%% LyX 2.4.4 created this file.  For more info, see https://www.lyx.org/.
%% Do not edit unless you really know what you are doing.
\documentclass[12pt,oneside,american]{amsart}
\usepackage[T1]{fontenc}
\usepackage[utf8]{inputenc}
\usepackage{amstext}
\usepackage{amsthm}
\usepackage{amssymb}
\usepackage{stackrel}
\usepackage{geometry}
\geometry{verbose,tmargin=2cm,bmargin=2cm,lmargin=2cm,rmargin=2cm}
\usepackage{wasysym}

\makeatletter
%%%%%%%%%%%%%%%%%%%%%%%%%%%%%% Textclass specific LaTeX commands.
\numberwithin{equation}{section}
\numberwithin{figure}{section}
\theoremstyle{plain}
\newtheorem{thm}{\protect\theoremname}[section]
\theoremstyle{definition}
\newtheorem{example}[thm]{\protect\examplename}
\theoremstyle{remark}
\newtheorem{rem}[thm]{\protect\remarkname}
\theoremstyle{definition}
\newtheorem{defn}[thm]{\protect\definitionname}
\theoremstyle{plain}
\newtheorem{lem}[thm]{\protect\lemmaname}
\theoremstyle{plain}
\newtheorem{cor}[thm]{\protect\corollaryname}
\theoremstyle{plain}
\newtheorem{prop}[thm]{\protect\propositionname}

%%%%%%%%%%%%%%%%%%%%%%%%%%%%%% User specified LaTeX commands.
\usepackage{enumitem}
\usepackage{amsmath}
\DeclareMathOperator{\tr}{tr}
\DeclareMathOperator{\dist}{dist}
\DeclareMathOperator{\Var}{Var}
\DeclareMathOperator{\ext}{Ext}
\DeclareMathOperator{\diam}{diam}
\DeclareMathOperator{\R}{\mathbb{R}}
\DeclareMathOperator{\id}{id}
\DeclareMathOperator{\C}{\mathbb{C}}
\DeclareMathOperator{\N}{\mathbb{N}}
\setenumerate{label=(\roman*)}
\DeclareMathOperator*{\argmin}{\arg\!\min}
\DeclareMathOperator*{\argmax}{\arg\!\max}
\allowdisplaybreaks
\usepackage[bb=dsserif]{mathalpha}
\usepackage{bm}
\newcommand{\ind}{\perp\!\!\!\!\perp} 

\makeatother

\usepackage{babel}
\providecommand{\corollaryname}{Corollary}
\providecommand{\definitionname}{Definition}
\providecommand{\examplename}{Example}
\providecommand{\lemmaname}{Lemma}
\providecommand{\propositionname}{Proposition}
\providecommand{\remarkname}{Remark}
\providecommand{\theoremname}{Theorem}

\begin{document}
\title{STAT 30100 Mathematical Statistics 1}
\author{Winter 2026}
\maketitle

\section{Sufficiency and Completeness}

\subsection{Sufficiency}

A statistical model is a collection of probability distributions $(P_{\theta}:\theta\in\Theta)$.
Then data or observations are random variables $X_{1},\ldots,X_{n}\sim P_{\theta}$
for some $\theta\in\Theta$. A statistic is a function of the data,
$T=T(X_{1},\ldots,X_{n}).$ A statistic $T$ is sufficient if the
conditional distribution $X|T$ does not depend on $\theta$. 
\begin{example}
If $X_{1},\ldots,X_{n}\overset{\text{iid}}{\sim}N(\theta,1)$ , then
$T(X)=\overline{X}$ is sufficient. We can see this by finding the
conditional distribution
\[
\left(\begin{array}{c}
X_{1}\\
\vdots\\
X_{n}
\end{array}\right)|\overline{X}\sim N\left(\left(\begin{array}{c}
\overline{X}\\
\vdots\\
\overline{X}
\end{array}\right),I_{n}-\frac{1}{n}\mathbb{{1}}_{n}\mathbb{{1}}_{n}^{T}\right).
\]
To see this, introduce $P:=\frac{1}{n}\mathbb{{1}}_{n}\mathbb{{1}}_{n}^{T}$
and write $X=PX+(I_{n}-P)X$. Then $PX=\overline{X}\mathbb{{1}}_{n}$,
and $(I_{n}-P)X\sim N(0,I_{n}-P)$. Thus, $X|\overline{X}\sim N\left(\overline{X}\mathbb{{1}}_{n},I_{n}-P\right).$
\end{example}

\begin{example}
If $X_{1},\ldots,X_{n}\overset{\text{iid}}{\sim}\text{Bernoulli}(\theta)$,
then $T(X)=\sum_{i=1}^{n}X_{i}$ is sufficient. This can be seen from
\begin{align*}
\mathbb{P}(X=x|T=t) & =\frac{\mathbb{P}(X=x,T=t)}{\mathbb{P}(T=t)}\\
 & =\frac{\mathbb{{1}}_{\{\sum_{i=1}^{n}x_{i}=t\}}\mathbb{P}(X=x)}{\mathbb{P}(T=t)}\\
 & =\frac{\mathbb{{1}}_{\{\sum_{i=1}^{n}x_{i}=t\}}\theta^{\sum_{i=1}^{n}x_{i}}(1-\theta)^{1-\sum_{i=1}^{n}x_{i}}}{\binom{n}{t}\theta^{t}(1-\theta)^{1-t}}\\
 & =\frac{\mathbb{{1}}_{\{\sum_{i=1}^{n}x_{i}=t\}}\theta^{t}(1-\theta)^{1-t}}{\binom{n}{t}\theta^{t}(1-\theta)^{1-t}}\\
 & =\frac{\mathbb{{1}}_{\{\sum_{i=1}^{n}x_{i}=t\}}}{\binom{n}{t}}.
\end{align*}
\end{example}

\begin{example}
If $X_{1},\ldots,X_{n}\overset{\text{iid}}{\sim}P_{\theta}$, then
$T=(X_{(1)},\ldots,X_{(n)})$ is sufficient because $X|T=t$ is the
uniform distribution over all $n!$ permutations of the vector $t$. 
\end{example}

\begin{example}
If $X_{1},\ldots,X_{n}\overset{\text{iid}}{\sim}\text{Unif}(0,\theta)$,
then $T=X_{(n)}$ is sufficient. This can be seen from
\begin{align*}
f_{X|T}(x|t) & =\frac{f_{X,T}(x,t)}{f_{T}(t)}\\
 & =\frac{\mathbb{{1}}_{\{\max(x_{1},\ldots,x_{n})=t\}}f_{X}(x)}{f_{T}(t)}\\
 & =\frac{\mathbb{{1}}_{\{\max(x_{1},\ldots,x_{n})=t\}}\left(\frac{1}{\theta}\right)^{n}}{n\left(\frac{t}{\theta}\right)^{n-1}\frac{1}{\theta}}\\
 & =\frac{\mathbb{{1}}_{\{\max(x_{1},\ldots,x_{n})=t\}}}{nt^{n-1}},
\end{align*}
where $f_{T}(t)=n\left(\frac{t}{\theta}\right)^{n-1}\frac{1}{\theta}$
follows from differentiating 
\[
F_{T}(t)=\mathbb{P}(T\leq t)=\left(\frac{t}{\theta}\right)^{n}.
\]
\end{example}

Shall we always use sufficient statistic and throw away the original
data? From the information-theoretic perspective, the answer is yes.
Indeed, if we can sample from $T$ and also from $X|(T=t)$ for an
arbitrary $t$, then we can sample from $X$. But from a computational
perspective, it is not always optimal because sampling from $X|T$
is often not easy or fast to do. In fact, it has been shown that sampling
$\widetilde{X}\sim X|T$ can be NP-hard. 
\begin{rem}
In a Bayesian setting where $\theta$ is random, another way to think
of sufficiency is that $\theta\to T\to X$ is a Markov chain.
\end{rem}

The following gives a characterization of sufficiency. 
\begin{thm}
Suppose that $(P_{\theta}:\theta\in\Theta)$ is continuous or discrete.
Then $T$ is sufficient if and only if $p(x|\theta)=g_{\theta}(T(x))h(x)$
for some measurable functions $g_{\theta}$ and $h$. 
\end{thm}

\begin{proof}
We will only prove this in the discrete case. Suppose that the pmf
can be factorized as $p(x|\theta)=g_{\theta}(T(x))h(x)$ for some
measurable functions $g_{\theta}$ and $h$. Then
\begin{align*}
\mathbb{P}(X=x,T=t) & =\mathbb{{1}}_{\{T(x)=t\}}\mathbb{P}(X=x)=\mathbb{{1}}_{\{T(x)=t\}}g_{\theta}(T(x))h(x)=\mathbb{{1}}_{\{T(x)=t\}}g_{\theta}(t)h(x),
\end{align*}
\[
\mathbb{P}(T=t)=\sum_{x':T(x')=t}\mathbb{P}(X=x')=\sum_{x':T(x')=t}g_{\theta}(T(x'))h(x')=g_{\theta}(t)\sum_{x':T(x')=t}h(x'),
\]
Therefore, 
\[
\mathbb{P}(X=x|T=t)=\frac{\mathbb{P}(X=x,T=t)}{\mathbb{P}(T=t)}=\frac{\mathbb{{1}}_{\{T(x)=t\}}h(x)}{\sum_{x':T(x')=t}h(x')}
\]
does not depend on $\theta$. Conversely, suppose that $T$ is sufficient.
Then
\begin{align*}
p(x|\theta) & =\mathbb{P}_{\theta}(X=x)\\
 & =\mathbb{P}_{\theta}(X=x,T(X)=T(x))\\
 & =\mathbb{P}(X=x|T(X)=T(x))\mathbb{P}_{\theta}(T(X)=T(x))\\
 & =h(x)g_{\theta}(T(x)),
\end{align*}
where $h(x)=\mathbb{P}(X=x|T(X)=T(x))$ and $g_{\theta}(T(x))=\mathbb{P}_{\theta}(T(X)=T(x))$. 
\end{proof}
%
\begin{example}
Let $X_{1},\ldots,X_{n}\overset{\text{iid}}{\sim}N(\theta,1).$ Then
\begin{align*}
p(x|\theta) & =\prod_{i=1}^{n}\frac{1}{\sqrt{2\pi}}e^{-\frac{1}{2}(x-\theta)^{2}}\\
 & =\left(\frac{1}{\sqrt{2\pi}}\right)^{n}e^{-\frac{1}{2}\sum_{i=1}^{n}(x_{i}-\theta)^{2}}\\
 & =\left(\frac{1}{\sqrt{2\pi}}\right)^{n}e^{-\frac{1}{2}\sum_{i=1}^{n}x_{i}^{2}-\frac{1}{2}n\theta^{2}+\theta\sum_{i=1}^{n}x_{i}}\\
 & =\left(\left(\frac{1}{\sqrt{2\pi}}\right)^{n}e^{-\frac{1}{2}\sum_{i=1}^{n}x_{i}^{2}}\right)e^{-\frac{1}{2}n\theta^{2}+\theta n\overline{x}}
\end{align*}
which shows that $T(X)=\overline{X}$ is sufficient. 
\end{example}

\begin{example}
If $X_{1},\ldots,X_{n}\overset{\text{iid}}{\sim}\text{Unif}(0,\theta)$,
then
\begin{align*}
p(x|\theta) & =\prod_{i=1}^{n}\left(\frac{1}{\theta}\mathbb{{1}}_{\{0<x_{i}<\theta\}}\right)=\theta^{-n}\mathbb{{1}}_{\{0<x_{(1)}\leq x_{(n)}<\theta\}}=\left(\mathbb{{1}}_{\{0<x_{(1)}\}}\right)\left(\theta^{-n}\mathbb{{1}}_{\{x_{(n)}<\theta\}}\right)
\end{align*}
which shows that $T(X)=X_{(n)}$ is sufficient. 
\end{example}


\subsection{Exponential families}
\begin{defn}
We say that a continuous or discrete statistical model $(P_{\theta}:\theta\in\Theta)$
is an exponential family if
\[
p(x|\theta)=\exp\left(\sum_{j=1}^{d}\eta_{j}(\theta)T_{j}(x)-B(\theta)\right)h(x)
\]
where $\eta(\theta)$ is called the natural parameter, $T(x)=(T_{1}(x),\ldots,T_{d}(x))^{T}$
is the sufficient statistic, 
\[
B(\theta)=\log\int\exp\left(\sum_{j=1}^{d}\eta_{j}(\theta)T_{j}(x)\right)h(x)\,d\mu(x)
\]
is the log-partition function, and $h\geq0$ is the base measure. 
\end{defn}

In the formula for $B(\theta)$ above, $\mu$ is either the Lebesgue
measure or the counting measure. Given observations $X_{1},\ldots,X_{n}\overset{\text{iid}}{\sim}p(x|\theta)=e^{\sum_{j=1}^{d}\eta_{j}(\theta)T_{j}(x)-B(\theta)}h(x)$,
we have
\[
p(x_{1},\ldots,x_{n}|\theta)=\exp\left(\sum_{j=1}^{d}\eta_{j}(\theta)\left(\sum_{i=1}^{n}T_{j}(x_{i})\right)-nB(\theta)\right)\prod_{i=1}^{n}h(x_{i}),
\]
from which we see that the joint distribution also is an exponential
family with the same natural parameter and sufficient statistic $T'(x_{1},\ldots,x_{n})=\sum_{i=1}^{n}T(x_{i})$.
We say that an exponential family is in canonical form if it is parameterized
by $\eta\in H$. That is,
\[
p(x|\eta)=\exp\left(\sum_{j=1}^{d}\eta_{j}T_{j}(x)-A(\eta)\right)h(x).
\]
Any exponential family can be reparameterized into canonical form.
Here, 
\[
A(\eta)=\log\int\exp\left(\sum_{j=1}^{d}\eta_{j}T_{j}(x)\right)h(x)\,d\mu(x).
\]

\begin{defn}
We say an exponential family $(P_{\eta}:\eta\in H)$ in canonical
form is minimal if the dimension $d$ cannot be reduced.
\end{defn}

\begin{example}
The family $p(x|\theta)=\exp(\eta_{1}T(x)+\eta_{2}(3T(x)+2)-A(\eta))$
is non-minimal, since 
\[
\exp(\eta_{1}T(x)+\eta_{2}(3T(x)+2)-A(\eta))=\exp((\eta_{1}+3\eta_{2})T(x)+2\eta_{2}-A(\eta)).
\]
Similarly, the family $p(x|\theta)=\exp\left(\eta T_{1}(x)+(4-5\eta)T_{2}(x)-A(\eta)\right)$
is non-minimal, since
\[
\exp\left(\eta T_{1}(x)+(4-5\eta)T_{2}(x)-A(\eta)\right)=\exp\left(\eta(T_{1}(x)-5T_{2}(x))-A(\eta)\right)\exp\left(4T_{2}(x)\right).
\]
\end{example}

Essentially, minimality fails if the components of the sufficient
statistic or the natural parameters are not linearly independent.
There are two types of minimal exponential families $(P_{\eta}:\eta\in H)$: 
\begin{enumerate}
\item full rank exponential families, where $H$ contains an open $d$-dimensional
rectangle;
\item curved exponential families, where $\eta_{1},\ldots,\eta_{d}$ are
related in a non-linear way. 
\end{enumerate}
%
\begin{example}
For $N(\mu,\sigma^{2})$, we have
\[
p(x|\mu,\sigma^{2})=\exp\left(-\frac{x^{2}}{2\sigma^{2}}+\frac{\mu}{\sigma^{2}}x-\frac{\mu^{2}}{2\sigma^{2}}-\frac{1}{2}\log(2\pi\sigma^{2})\right),
\]
so $T_{1}(x)=-x^{2}$, $T_{2}(x)=x$, $\eta_{1}=\frac{1}{2\sigma^{2}}$,
$\eta_{2}=\frac{\mu}{\sigma^{2}}$. Let's consider different special
cases. 
\begin{enumerate}
\item If $\mu=\sigma^{2}$, then $\eta_{2}=1$, so this is a non-minimal
exponential family. 
\item If $\mu=\sqrt{\sigma^{2}}$, then $\eta_{1}=\frac{1}{2\sigma^{2}}$
and $\eta_{2}=\frac{1}{\sqrt{\sigma^{2}}}$, so $\eta^{2}=2\eta_{1}$.
This is a minimal curved exponential family. 
\item If $\mu$ and $\sigma^{2}$ do not have constraints, then this is
a full rank minimal exponential family. 
\end{enumerate}
\end{example}


\subsection{Minimal sufficiency and completeness}

We say that a statistic $S$ is minimal sufficient if it is sufficient
and for every sufficient $T$, there is a measurable function $f$
such that $S=f(T)$. 

\subsubsection{Sub-family method}

How do we find minimal sufficient statistic? The first method is called
the sub-family method. 
\begin{lem}
Suppose that $\Theta_{0}\subseteq\Theta$ and $S$ is minimal sufficient
for $(P_{\theta}:\theta\in\Theta_{0})$ and sufficient for $(P_{\theta}:\theta\in\Theta)$.
If the support of $P_{\theta}$ does not depend on $\theta$, then
$S$ is minimal sufficient for $(P_{\theta}:\theta\in\Theta)$. 
\end{lem}

\begin{proof}
Suppose that $T$ is sufficient for $(P_{\theta}:\theta\in\Theta)$.
Then $T$ is sufficient for $(P_{\theta}:\theta\in\Theta_{0})$, so
since $S$ is minimal sufficient for $(P_{\theta}:\theta\in\Theta_{0})$,
it follows that $S$ is a function of $T$. 
\end{proof}
%
\begin{thm}
For a discrete or continuous statistical model $(P_{\theta}:\theta\in\{\theta_{0},\ldots,\theta_{d}\})$
with finite parameter space and common support, 
\[
T(X)=\left(\frac{p(X|\theta_{1})}{p(X|\theta_{0})},\frac{p(X|\theta_{2})}{p(X|\theta_{0})},\ldots,\frac{p(X|\theta_{d})}{p(X|\theta_{0})}\right)
\]
is a minimal sufficient statistic for $(P_{\theta}:\theta\in\{\theta_{0},\ldots,\theta_{d}\})$. 
\end{thm}

\begin{proof}
For $j=1,\ldots,d$, we have $p(x|\theta_{j})=T_{j}(x)p(x|\theta_{0})$,
so if we define
\[
g_{\theta_{j}}(T(x))=\begin{cases}
1 & \text{if }j=0,\\
T_{j}(x) & \text{if }j=1,\ldots,d,
\end{cases}\quad\text{and}\quad h(x)=p(x|\theta_{0}),
\]
we see that $T$ is sufficient by the factorization theorem. For any
sufficient $T'$, the factorization theorem implies the likelihood
ratio is a function of $T'$, so $T$ is a function of $T'$. 
\end{proof}
\begin{example}
Let $X_{1},\ldots,X_{n}\overset{\text{iid}}{\sim}P_{\theta}$ where
$P_{\theta}=\text{Bernoulli}(\theta)$ and $\theta\in(0,1)$. Consider
$\theta_{0}=0.3$ and $\theta_{1}=0.6$. Then
\[
\frac{p(X|\theta_{1})}{p(X|\theta_{0})}=\frac{\theta_{1}^{\sum_{i=1}^{n}X_{i}}(1-\theta_{1})^{n-\sum_{i=1}^{n}X_{i}}}{\theta_{0}^{\sum_{i=1}^{n}X_{i}}(1-\theta_{0})^{n-\sum_{i=1}^{n}X_{i}}}=\left(\frac{\theta_{1}(1-\theta_{0})}{\theta_{0}(1-\theta_{1})}\right)^{\sum_{i=1}^{n}X_{i}}\left(\frac{1-\theta_{1}}{1-\theta_{0}}\right)^{n}=2^{\sum_{i=1}^{n}X_{i}}\left(\frac{4}{7}\right)^{n}.
\]
This statistic is equivalent to $T(X)=\sum_{i=1}^{n}X_{i}$, so by
the previous theorem we see that $T(X)$ is minimal sufficient for
the sub-family $(P_{\theta}:\theta\in\{0.3,0.6\})$. We have seen
before that $T(X)$ is sufficient for $(P_{\theta}:\theta\in[0,1])$,
so the sub-family lemma yields that $T(X)$ is minimal sufficient
for $(P_{\theta}:\theta\in[0,1])$. 
\end{example}

\begin{thm}
If $(P_{\theta}:\theta\in H)$ is a minimal exponential family in
canonical form with
\[
p(x|\eta)=\exp\left(\eta^{T}T(x)-A(\eta)\right)h(x),\quad\eta\in H,
\]
then $T(X)$ is minimal sufficient. 
\end{thm}

\begin{proof}
Since the exponential family is minimal, we can find $\eta_{0},\ldots,\eta_{d}\in H$
such that
\[
\left(\begin{array}{c}
(\eta_{1}-\eta_{0})^{T}\\
(\eta_{2}-\eta_{1})^{T}\\
\vdots\\
(\eta_{d}-\eta_{0})^{T}
\end{array}\right)\in\mathbb{R}^{d\times d}
\]
has full rank. To see this, note that if the exponential family is
full rank, then $H$ is an open rectangle, so it immediately follows.
If the exponential family is curved, then this follows from the non-zero
curvature of the function which relates the components of $\eta$.
Then we know that
\[
\left(\frac{p(X|\eta_{1})}{p(X|\eta_{0})},\ldots,\frac{p(X|\eta_{d})}{p(X|\eta_{0})}\right)
\]
is minimal sufficient for the $(P_{\eta}:\eta\in\{\eta_{1},\ldots,\eta_{d}\})$.
Note that
\[
\frac{p(X|\eta_{j})}{p(X|\eta_{0})}=\frac{\exp\left(\langle\eta_{j},T(X)\rangle-A(\eta_{j})\right)}{\exp\left(\langle\eta_{0},T(X)\rangle-A(\eta_{0})\right)}=\exp\left(\left\langle \eta_{j}-\eta_{0},T(X)\right\rangle -A(\eta_{j})+A(\eta_{0})\right),
\]
which is equivalent to
\[
\left(\begin{array}{c}
\left\langle \eta_{1}-\eta_{0},T(x)\right\rangle \\
\left\langle \eta_{2}-\eta_{0},T(x)\right\rangle \\
\vdots\\
\left\langle \eta_{d}-\eta_{0},T(x)\right\rangle 
\end{array}\right)=\left(\begin{array}{c}
(\eta_{1}-\eta_{0})^{T}\\
(\eta_{2}-\eta_{0})^{T}\\
\vdots\\
(\eta_{d}-\eta_{0})^{T}
\end{array}\right)T(X),
\]
which is equivalent to $T(X)$. Since $T(X)$ is sufficient for the
whole family $(P_{\eta}:\eta\in H)$, it follows from the sub-family
lemma that $T(X)$ is minimal sufficient. 
\end{proof}
%
One of the restrictions of the sub-family method is that the support
of all the distributions in the family must have the same support.
This is not an issue for exponential families as the support depends
only on the base measure $h(x)$, but it could be an issue for other
distributions like $P_{\theta}=\text{Unif}(0,\theta)$. 

\subsubsection{Completeness}

The second method of finding minimal sufficient statistics is called
completeness. The idea is to remove ancillary information. For example,
if $X_{1}$ and $X_{2}$ are i.i.d. from $(N(\theta,1):\theta\in\mathbb{R})$,
then we know that $(X_{1},X_{2})$ is not a minimal sufficient statistic
because it is equivalent to $(X_{1}-X_{2},X_{1}+X_{2})$, and $X_{1}-X_{2}\sim N(0,2)$
carries no information about $\theta$. 
\begin{defn}
We say that a statistic $A=A(X)$ is ancillary if its distribution
does not depend on $\theta$, and we say it is first-order ancillary
if $\mathbb{E}_{\theta}[A(X)]$ does not depend on $\theta$. 
\end{defn}

\begin{defn}
We say a statistic $T=T(X)$ is complete if $\mathbb{E}_{\theta}[f(T(X))]=0$
for all $\theta\in\Theta$ implies $\mathbb{P}_{\theta}[f(T(X))=0]=1$
for all $\theta\in\Theta$. In other words, every function of $T$
which is first-order ancillary must be constant almost surely. 
\end{defn}

\begin{thm}
[Bahadur] If $T$ is sufficient and complete, then it is minimal
sufficient.
\end{thm}

\begin{proof}
[Sketch of proof] Assume that a minimal sufficient statistic $U$
exists. By definition, we have $U=h(T)$ for some measurable function
$h$. We want to show that $T$ is also a function of $U$. Define
$g(u)=\mathbb{E}_{\theta}[T|U=u]$. Note that $g$ does not depend
on $\theta$ because $U$ is sufficient. Then $\mathbb{E}_{\theta}[g(h(T))]=\mathbb{E}_{\theta}[g(U)]=\mathbb{E}_{\theta}\left[\mathbb{E}_{\theta}\left[T|U\right]\right]=\mathbb{E}_{\theta}[T]$
, so $\mathbb{E}_{\theta}\left[g(h(T))-T\right]=0$ for all $\theta\in\Theta$.
By completeness of $T$, it follows that $g(U)=g(h(T))=T$ almost
surely for all $\theta\in\Theta$. 
\end{proof}
%
\begin{example}
Let $X_{1},\ldots,X_{n}\overset{\text{iid}}{\sim}\text{Bernoulli}(\theta)$
where $\theta\in(0,1)$. Let $T(X)=\sum_{i=1}^{n}X_{i}\sim\text{Binomial}(n,\theta)$.
Suppose $\mathbb{E}_{\theta}f(T(X))=0$ for all $\theta\in(0,1)$.
Then 
\[
\mathbb{E}_{\theta}[f(T(X))]=\sum_{i=1}^{n}f(i)\binom{n}{i}\theta^{i}(1-\theta)^{n-i}=(1-\theta)^{n}\sum_{i=1}^{n}f(i)\binom{n}{i}\left(\frac{\theta}{1-\theta}\right)^{i},
\]
so we have
\[
\sum_{i=1}^{n}f(i)\binom{n}{i}\left(\frac{\theta}{1-\theta}\right)^{i}=0\quad\text{for all }\theta\in(0,1).
\]
Set $\beta=\frac{\theta}{1-\theta}$. Then 
\[
\sum_{i=1}^{n}f(i)\binom{n}{i}\beta^{i}=0\quad\text{for all }\beta>0.
\]
This means that the polynomial $p(x)=\sum_{i=1}^{n}f(i)\binom{n}{i}x^{i}=0$
has infinitely many roots, which can only occur if all coefficients
vanish. Therefore, $f(i)=0$ for all $i=1,\ldots,n$. This shows that
$T$ is complete. We have seen before that $T(X)$ is sufficient,
so Bahadur's theorem implies it is also minimal sufficient. 
\end{example}

\begin{example}
Let $X_{1},\ldots,X_{n}\overset{\text{iid}}{\sim}\text{Unif}(0,\theta)$
where $\theta>0$. Let $T(X)=X_{(n)}$. Suppose $\mathbb{E}_{\theta}f(T(X))=0$
for all $\theta>0$. Then 
\begin{align*}
\theta^{-n}n\int_{0}^{\theta}t^{n-1}f(t)\,dt=0\quad\text{for all }\theta>0 & .
\end{align*}
This implies that $G(\theta)=\int_{0}^{\theta}t^{n-1}f(t)\,dt=0$
for all $\theta>0$. Then by the Lebesgue differentiation theorem,
$G'(\theta)=\theta^{n-1}f(\theta)$ for almost every $\theta>0$.
Since $G=0$, we have $G'=0$, so $f=0$ almost everywhere on $(0,\infty)$.
This implies that $\mathbb{P}_{\theta}[f(T(X))=0]=1$ for all $\theta\in\Theta$. 
\end{example}

\begin{example}
Let $X_{1},\ldots,X_{n}\overset{\text{iid}}{\sim}N(\theta,1)$ where
$\theta\in\mathbb{R}$. Let $T(X)=\frac{1}{\sqrt{n}}\sum_{i=1}^{n}X_{i}\sim N(\sqrt{n}\theta,1)=N(\beta,1)$,
where $\beta=\sqrt{n}\theta$. Note that
\[
\mathbb{E}_{\theta}[f(T(X))]=\int_{\mathbb{R}}f(x)e^{-\frac{1}{2}(x-\beta)^{2}}\,dx.
\]
Suppose $\mathbb{E}_{\theta}[f(T(X))]=0$ for all $\theta\in\mathbb{R}$,
so that $\int_{\mathbb{R}}f(x)e^{-\frac{1}{2}x^{2}+x\beta}\,dx=0$
for all $\beta\in\mathbb{R}$. Decomposing $f=f^{+}-f^{-}$, we have
\[
\int_{\mathbb{R}}f^{+}(x)e^{-\frac{1}{2}x^{2}+x\beta}\,dx=\int_{\mathbb{R}}f^{-}(x)e^{-\frac{1}{2}x^{2}+x\beta}\,dx\quad\text{for all }\beta\in\mathbb{R}.
\]
Then 
\[
\frac{\int_{\mathbb{R}}f^{+}(x)e^{-\frac{1}{2}x^{2}}e^{x\beta}\,dx}{\int_{\mathbb{R}}f^{+}(t)e^{-\frac{1}{2}t^{2}}\,dt}=\frac{\int_{\mathbb{R}}f^{-}(x)e^{-\frac{1}{2}x^{2}}e^{x\beta}\,dx}{\int_{\mathbb{R}}f^{-}(t)e^{-\frac{1}{2}t^{2}}\,dt}\quad\text{for all }\beta\in\mathbb{R}.
\]
This shows that the MGFs of the densities $\frac{f^{+}(x)e^{-\frac{1}{2}x^{2}}}{\int_{\mathbb{R}}f^{+}(t)e^{-\frac{1}{2}t^{2}}\,dt}$
and $\frac{f^{-}(x)e^{-\frac{1}{2}x^{2}}}{\int_{\mathbb{R}}f^{-}(t)e^{-\frac{1}{2}t^{2}}\,dt}$
coincide, so $f^{+}(x)e^{-\frac{1}{2}x^{2}}=f^{-}(x)e^{-\frac{1}{2}x^{2}}$
for all $x\in\mathbb{R}$ and we must have $f^{+}=f^{-}$, so $f=0$.
Thus, $T(X)$ is complete, which implies that $\overline{X}$ is sufficient
and complete and hence minimal sufficient. 
\end{example}

\begin{thm}
For full rank minimal exponential family $(P_{\eta}:\eta\in H)$,
$T$ is complete.
\end{thm}

\begin{thm}
[Basu] Let $T$ be complete and sufficient. If $A$ is ancillary,
then $T\ind A$. 
\end{thm}

\begin{proof}
Let $B$ be a measurable set. We want to show that $\mathbb{P}_{\theta}[A\in B|T]=\mathbb{P}_{\theta}[A\in B]$
almost surely for all $\theta\in\Theta$. Set $c=\mathbb{P}_{\theta}(A\in B)$,
which does not depend on $\theta$ since $A$ is ancillary. Also,
set $g(T)=\mathbb{P}_{\theta}(A\in B|T)$, which also does not depend
on $\theta$ since $T$ is sufficient. Then 
\[
\mathbb{E}_{\theta}[g(T)-c]=\mathbb{E}_{\theta}[\mathbb{P}_{\theta}[A\in B|T]]-\mathbb{P}_{\theta}[A\in B]=0\quad\text{for all }\theta\in\Theta,
\]
so by completeness we have $g(T)=c$ almost surely for all $\theta\in\Theta$. 
\end{proof}
%
\begin{example}
Let $X_{1},\ldots,X_{n}\overset{\text{iid}}{\sim}N(\theta,1)$. Then
$\overline{X}\ind\sum_{i=1}^{n}(X_{i}-\overline{X})^{2}$ by Basu's
theorem since $\sum_{i=1}^{n}(X_{i}-\overline{X})^{2}\sim\chi_{n-1}^{2}$
is ancillary. 
\end{example}


\section{Decision Theory}

\subsection{Bayes and minimax estimators}

Let $X_{1},\ldots,X_{n}\overset{\text{iid}}{\sim}P_{\theta}$, and
let $\widehat{\theta}=\widehat{\theta}(X_{1},\ldots,X_{n})$ be an
estimator of the parameter $\theta$. Let $L(\widehat{\theta},\theta)$
be a loss function which quantifies the accuracy of the estimator.
Then the risk function is
\[
R(\widehat{\theta},\theta)=\mathbb{E}_{\theta}[L(\widehat{\theta},\theta)]=\int L(\widehat{\theta}(x),\theta)p(x|\theta)\,dx
\]

\begin{thm}
[Rao-Blackwell] Assume $L(\widehat{\theta},\theta)$ is convex in
$\widehat{\theta}$. For any $\widehat{\theta}$ and any sufficient
$T$, define $\widetilde{\theta}=\mathbb{E}_{\theta}(\widehat{\theta}|T)$.
Then $R(\widetilde{\theta},\theta)\leq R(\widehat{\theta},\theta).$
\end{thm}

\begin{proof}
We have $L(\widetilde{\theta},\theta)=L(\mathbb{E}_{\theta}(\widehat{\theta}|T),\theta)\leq\mathbb{E}_{\theta}(L(\widehat{\theta},\theta)|T)$,
and taking expectation of both sides gives $R(\widetilde{\theta},\theta)\leq R(\widehat{\theta},\theta).$
\end{proof}
%
\begin{defn}
Given an estimator $\widehat{\theta}$, we define the maximum risk
to be $\sup_{\theta\in\Theta}R(\widehat{\theta},\theta)$. We say
that $\widehat{\theta}$ is minimax if $\widehat{\theta}\in\argmin_{\widetilde{\theta}}\sup_{\theta\in\Theta}R(\widetilde{\theta},\theta)$
where the minimum is taken over all estimators. Given an estimator
$\widehat{\theta}$ and a prior distribution $\pi$ on $\Theta$,
we define the average risk with respect to $\pi$ to be $\int R(\widehat{\theta},\theta)\pi(\theta)\,d\theta$.
We say that $\widehat{\theta}$ is a Bayes estimator if $\widehat{\theta}\in\argmin_{\widetilde{\theta}}\int_{\Theta}R(\widetilde{\theta},\theta)\pi(\theta)\,d\theta$. 
\end{defn}


\subsubsection{Finding Bayes estimators}
\begin{lem}
The Bayes estimator is given by 
\[
\widehat{\theta}_{\pi}(X)=\argmin_{a}\int L(a,\theta)\pi(\theta|X)\,d\theta.
\]
\end{lem}

\begin{proof}
Let $m(x)=\int p(x|\theta)\pi(\theta)\,d\theta$ be the marginal of
$X$. For any estimator $\widetilde{\theta}$, we have
\begin{align*}
\int R(\widehat{\theta}_{\pi},\theta)\pi(\theta)\,d\theta & =\int\int L(\widehat{\theta}_{\pi}(x),\theta)p(x|\theta)\pi(\theta)\,dx\,d\theta\\
 & =\int\int L(\widehat{\theta}_{\pi}(x),\theta)\pi(\theta|x)d\theta\,m(x)\,dx\\
 & \leq\int\int L(\widetilde{\theta}(x),\theta)\pi(\theta|x)d\theta\,m(x)\,dxx\\
 & =\int R(\widetilde{\theta},\theta)\pi(\theta)\,d\theta.
\end{align*}
\end{proof}
%
\begin{cor}
If $\Theta\subseteq\mathbb{R}$ and $L(\widehat{\theta},\theta)=(\widehat{\theta}-\theta)^{2}$.
Then 
\[
\widehat{\theta}_{\pi}(X)=\argmin_{a}\int(a-\theta)^{2}\pi(\theta|X)\,d\theta=\argmin_{a}\mathbb{E}[(a-\theta)^{2}|X]=\mathbb{E}[\theta|X].
\]
\end{cor}

\begin{example}
Let $X_{1},\ldots,X_{n}\overset{\text{i.i.d.}}{\sim}\text{Bernoulli}(\theta)$,
and let $L(\hat{\theta},\theta)=(\hat{\theta}-\theta)^{2}$. Consider
the prior $\pi\sim\text{Beta}(\alpha,\beta)$, so $\pi(\theta)\propto\theta^{\alpha-1}(1-\theta)^{\beta-1}$.
Then
\[
\pi(\theta|X)\propto\left(\prod_{i=1}^{n}\theta^{x_{i}}(1-\theta)^{1-x_{i}}\right)\pi(\theta)\propto\theta^{\alpha+\sum_{i=1}^{n}X_{i}-1}(1-\theta)^{n-\sum_{i=1}^{n}X_{i}+\beta-1}.
\]
This shows that the posterior distribution is $\theta|X_{1},\ldots,X_{n}\sim\text{Beta}\left(\sum_{i=1}^{n}X_{i}+\alpha,n-\sum_{i=1}^{n}X_{i}+\beta\right).$
The Bayes estimator is then
\[
\widehat{\theta}(X)=\mathbb{E}[\theta|X_{1},\ldots,X_{n}]=\frac{\sum_{i=1}^{n}X_{i}+\alpha}{n+\alpha+\beta}.
\]
\end{example}


\subsubsection{Finding minimax estimators}
\begin{thm}
If there exists a prior $\pi$ on $\Theta$ such that $\widehat{\theta}$
satisfies 
\[
\sup_{\theta\in\Theta}R(\widehat{\theta},\theta)=\inf_{\widetilde{\theta}}\int R(\widetilde{\theta},\theta)\pi(\theta)\,d\theta,
\]
then $\widehat{\theta}$ is minimax. 
\end{thm}

\begin{proof}
For all $\widetilde{\theta}$, we have
\[
\sup_{\theta\in\Theta}R(\widetilde{\theta},\theta)\geq\int R(\widetilde{\theta},\theta)\pi(\theta)\,d\theta\geq\inf_{\widetilde{\theta}}\int R(\widetilde{\theta},\theta)\pi(\theta)\,d\theta=\sup_{\theta\in\Theta}R(\widehat{\theta},\theta).
\]
\end{proof}
%
For a prior satisfying the above condition, we call it a least-favorable
prior. 
\begin{cor}
If $\widehat{\theta}$ is Bayes with respect to some prior $\pi$
and the risk $R(\widehat{\theta},\theta)$ is constant over $\theta\in\Theta$,
then $\widehat{\theta}$ is minimax. 
\end{cor}

\begin{proof}
We have $\sup_{\theta\in\Theta}R(\widehat{\theta},\theta)=\int R(\widehat{\theta},\theta)\pi(\theta)\,d\theta=\inf_{\widetilde{\theta}}\int R(\widetilde{\theta},\theta)\pi(\theta)\,d\theta$,
so by the theorem we have that $\widehat{\theta}$ is minimax. 
\end{proof}
%
The above corollary is usually the first line of attack when finding
minimax estimators.
\begin{example}
Let $X_{1},\ldots,X_{n}\overset{\text{iid}}{\sim}\text{Bernoulli}(\theta)$,
and let $L(\widehat{\theta},\theta)=(\hat{\theta}-\theta)^{2}$. Recall
that for the prior $\pi\sim\text{Beta}(\alpha,\beta)$ we have the
Bayes estimator
\[
\widehat{\theta}=\mathbb{E}[\theta|X_{1},\ldots,X_{n}]=\frac{\sum_{i=1}^{n}X_{i}+\alpha}{\alpha+\beta+n}=\frac{\alpha}{\alpha+\beta+n}+\frac{n}{\alpha+\beta+n}\overline{X}.
\]
The risk is 
\begin{align*}
R(\widehat{\theta},\theta) & =\mathbb{E}_{\theta}[(\widehat{\theta}-\theta)^{2}]\\
 & =\left(\mathbb{E}_{\theta}[\widehat{\theta}]-\theta\right)^{2}+\operatorname{Var}[\widehat{\theta}]\\
 & =\left(\frac{\alpha}{\alpha+\beta+n}+\frac{n}{\alpha+\beta+n}\theta-\theta\right)^{2}+\left(\frac{n}{\alpha+\beta+n}\right)^{2}\frac{\theta(1-\theta)}{n}\\
 & =\left(\frac{\alpha}{\alpha+\beta+n}-\frac{\alpha+\beta}{\alpha+\beta+n}\theta\right)^{2}+\frac{n}{(\alpha+\beta+n)^{2}}\theta(1-\theta)\\
 & =\left[\left(\frac{\alpha+\beta}{\alpha+\beta+n}\right)^{2}-\frac{n}{\alpha+\beta+n}\right]\theta^{2}+\left[\frac{n}{(\alpha+\beta+n)^{2}}-\frac{2\alpha(\alpha+\beta)}{\left(\alpha+\beta+n\right)^{2}}\right]\theta+\left(\frac{\alpha+\beta}{\alpha+\beta+n}\right)^{2}.
\end{align*}
If $\alpha=\beta=\frac{\sqrt{n}}{2}$, then the coefficients of the
$\theta$ and $\theta^{2}$ terms vanish, and the risk becomes $R(\widehat{\theta},\theta)=\frac{1}{4(\sqrt{n}+1)^{2}}$
which is constant with respect to $\theta$. Therefore, a minimax
estimator is
\[
\widehat{\theta}_{\text{minimax}}=\frac{\sum_{i=1}^{n}X_{i}+\frac{\sqrt{n}}{2}}{n+\sqrt{n}}.
\]
How does this estimator compare to the maximum likelihood estimator
$\widehat{\theta}_{\text{MLE}}=\overline{X}$? The risk function for
the MLE estimator is 
\[
R(\widehat{\theta}_{\text{MLE}},\theta)=\mathbb{E}_{\theta}[(\widehat{\theta}_{\text{MLE}}-\theta)^{2}]=\operatorname{Var}[\widehat{\theta}_{\text{MLE}}]=\frac{\theta(1-\theta)}{n}.
\]
There are some situations where the MLE has lower risk, for example
when $\theta$ is close to 0 or 1. However, if we consider the worst
case, we find that
\[
\sup_{\theta\in(0,1)}R(\widehat{\theta}_{\text{MLE}},\theta)=\frac{1}{4n}>\frac{1}{4(\sqrt{n}+1)^{2}}=\sup_{\theta\in(0,1)}R(\hat{\theta}_{\text{minimax}},\theta).
\]
\end{example}

The choice of loss function is important. Let us demonstrate this
with another example. 
\begin{example}
Let $X_{1},\ldots,X_{n}\overset{\text{iid}}{\sim}\text{Bernoulli}(\theta)$,
and let $L(\widehat{\theta},\theta)=\frac{(\widehat{\theta}-\theta)^{2}}{\theta(1-\theta)}$.
For the prior $\pi(\theta)=1$, the Bayes estimator is
\begin{align*}
\widehat{\theta}(X) & =\argmin_{a}\int\frac{(a-\theta)^{2}}{\theta(1-\theta)}\pi(\theta|X)\,d\theta=\argmin_{a}\int(a-\theta)^{2}\frac{\pi(\theta|X)}{\theta(1-\theta)}\,d\theta,
\end{align*}
where 
\[
\frac{\pi(\theta|X)}{\theta(1-\theta)}\propto\theta^{\sum_{i=1}^{n}X_{i}-1}(1-\theta)^{n-\sum_{i=1}^{n}X_{i}-1}=\text{Beta}\left(\sum_{i=1}^{n}X_{i},n-\sum_{i=1}^{n}X_{i}\right),
\]
so
\[
\widehat{\theta}(X)=\frac{\sum_{i=1}^{n}X_{i}}{\sum_{i=1}^{n}X_{i}+n-\sum_{i=1}^{n}X_{i}}=\overline{X}=\widehat{\theta}_{\text{MLE}}.
\]
The risk is $R(\widehat{\theta},\theta)=\frac{\mathbb{E}_{\theta}(\widehat{\theta}-\theta)^{2}}{\theta(1-\theta)}=\frac{1}{n}$
which is constant, so $\overline{X}$ is a Bayes estimator (which
cannot be the case with the squared loss function because with this
choice of loss only biased estimators could be Bayes), and it is also
a minimax estimator. 
\end{example}

Sometimes we may not be able to show that $\widehat{\theta}$ is a
Bayes estimator. In this case, we can find a sequence of Bayes estimators
whose Bayes risks approximate the maximum risk of $\widehat{\theta}$.
\begin{thm}
If there exists a sequence of priors $(\pi_{m})$ such that 
\[
\sup_{\theta\in\Theta}R(\widehat{\theta},\theta)=\lim_{m\to\infty}\inf_{\widetilde{\theta}}\int R(\widetilde{\theta},\theta)\pi_{m}(\theta)\,d\theta,
\]
then $\widehat{\theta}$ is minimax. 
\end{thm}

\begin{proof}
For all $\widetilde{\theta}$ and $m$, we have
\[
\sup_{\theta\in\Theta}R(\widetilde{\theta},\theta)\geq\int R(\widetilde{\theta},\theta)\pi_{m}(\theta)\,d\theta\geq\inf_{\widetilde{\theta}}\int R(\widetilde{\theta},\theta)\pi_{m}(\theta)\,d\theta.
\]
Taking the limit as $m\to\infty$ gives $\sup_{\theta\in\Theta}R(\widetilde{\theta},\theta)\geq\sup_{\theta\in\Theta}R(\widehat{\theta},\theta)$. 
\end{proof}
\begin{example}
Let $X_{1},\ldots,X_{n}\overset{\text{i.i.d.}}{\sim}N(\theta,1)$
where $\theta\in\mathbb{R}$. Consider the loss function $L(\widehat{\theta},\theta)=(\widehat{\theta}-\theta)^{2}$.
We claim that the estimator $\overline{X}$ is minimax. This is true
intuitively because the risk $R(\overline{X},\mu)=\mathbb{E}_{\theta}[(\overline{X}-\theta)^{2}]=\frac{1}{n}$
is constant. To prove this formally, consider the prior $\pi\sim N(0,\tau^{2})$.
Then
\[
\pi(\theta|X_{1},\ldots,X_{n})\propto\exp\left(-\frac{1}{2}\sum_{i=1}^{n}(X_{i}-\theta)^{2}-\frac{\theta^{2}}{2\tau^{2}}\right).
\]
This is a normal distribution, so the posterior mean is equal to the
posterior mode. Let $f(\theta):=\frac{1}{2}\sum_{i=1}^{n}(X_{i}-\theta)^{2}+\frac{\theta^{2}}{2\tau^{2}}$.
Then $f'(\theta)=\sum_{i=1}^{n}(\theta-X_{i})+\frac{\theta}{\tau^{2}}$,
and setting this to 0 gives the Bayes estimator
\[
\widehat{\theta}_{\tau^{2}}=\frac{n}{n+\tau^{-2}}\overline{X}.
\]
The risk is
\begin{align*}
R(\hat{\theta}_{\tau^{2}},\theta) & =\mathbb{E}[(\hat{\theta}_{\tau^{2}}-\theta)^{2}]\\
 & =\left(\frac{n}{n+\tau^{-2}}\theta-\theta\right)^{2}+\left(\frac{n}{n+\tau^{-2}}\right)^{2}\frac{1}{n}\\
 & =\left(\frac{\tau^{2}}{n+\tau^{-2}}\right)^{2}\theta^{2}+\left(\frac{n}{n+\tau^{-2}}\right)^{2}\frac{1}{n}.
\end{align*}
The average risk is
\[
\int R(\hat{\theta}_{\tau^{2}},\theta)\pi_{\tau^{2}}(\theta)\,d\theta=\left(\frac{\tau^{-2}}{n+\tau^{-2}}\right)^{2}\tau^{2}+\left(\frac{n}{n+\tau^{-2}}\right)^{2}\frac{1}{n}=\frac{1}{n+\tau^{-2}}.
\]
Thus, we have
\[
\sup_{\theta\in\mathbb{R}}R(\overline{X},\theta)=\lim_{\tau\to\infty}\int R(\hat{\theta}_{\tau^{2}},\theta)\pi_{\tau^{2}}(\theta)\,d\theta
\]
which shows that $\overline{X}$ is minimax. 
\end{example}


\subsection{Admissibility}

An estimator $\widehat{\theta}$ is inadmissible if there exists another
estimator $\widetilde{\theta}$ such that 
\begin{enumerate}
\item $R(\widetilde{\theta},\theta)\leq R(\widehat{\theta},\theta)$ for
all $\theta\in\Theta$, and
\item $R(\widetilde{\theta},\theta_{0})<R(\widehat{\theta},\theta_{0})$
for some $\theta_{0}\in\Theta$.
\end{enumerate}
Then we say $\widehat{\theta}$ is admissible if it is not inadmissible.
Inadmissible estimators are those that are essentially useless. There
are always other choices that are strictly better. 
\begin{thm}
Assume that the loss function is continuous. If $\widehat{\theta}$
is Bayes with respect to a prior $\pi$ whose support is the entire
parameter space $\Theta$, then it is admissible. 
\end{thm}

\begin{proof}
Suppose that $\widehat{\theta}$ is inadmissible. Then there exists
$\widetilde{\theta}$ such that $R(\widetilde{\theta},\theta)\leq R(\widehat{\theta},\theta)$
for all $\theta\in\Theta$, and $R(\widetilde{\theta},\theta_{0})<R(\widehat{\theta},\theta_{0})$
for some $\theta_{0}\in\Theta$. Then there exists an open set $\Theta_{0}\ni\theta_{0}$
and $\epsilon>0$ such that $R(\widetilde{\theta},\theta)<R(\widehat{\theta},\theta)-\varepsilon$
for all $\theta\in\Theta_{0}$. Then
\begin{align*}
\int_{\Theta}R(\widetilde{\theta},\theta)\pi(\theta)\,d\theta & =\int_{\Theta_{0}}R(\widetilde{\theta},\theta)\pi(\theta)\,d\theta+\int_{\Theta\backslash\Theta_{0}}R(\widetilde{\theta},\theta)\pi(\theta)\,d\theta\\
 & \leq\int_{\Theta_{0}}R(\widehat{\theta},\theta)\pi(\theta)\,d\theta+\int_{\Theta\backslash\Theta_{0}}R(\widehat{\theta},\theta)\pi(\theta)\,d\theta-\varepsilon\pi(\Theta_{0})\\
 & <\int_{\Theta}R(\widehat{\theta},\theta)\pi(\theta)\,d\theta,
\end{align*}
so $\widehat{\theta}$ cannot be Bayes. 
\end{proof}
%
Is the converse true? That is, is it true that an admissible estimator
must be Bayes? This is answered by the complete class theorem, which
roughly states that if $\widehat{\theta}$ is admissible, then there
exists a sequence of priors $(\pi_{m})$ such that $\widehat{\theta}_{\pi_{m}}\to\widehat{\theta}$. 
\begin{thm}
Let $X_{1},\ldots,X_{n}\overset{\mathrm{iid}}{\sim}N(\theta,1)$,
$\theta\in\mathbb{R}$. Consider the loss function $L(\hat{\theta},\theta)=(\hat{\theta}-\theta)^{2}$.
Then $\overline{X}$ is admissible.
\end{thm}

\begin{proof}
The proof is by contradiction. Suppose that $\widehat{\theta}=\overline{X}$
is inadmissible, so that there exists $\widetilde{\theta}$ such that
$R(\widetilde{\theta},\theta)\leq\frac{1}{n}$ for all $\theta\in\mathbb{R}$,
and $R(\widetilde{\theta},\theta_{0})<\frac{1}{n}$ for some $\theta_{0}\in\mathbb{R}$.
Then there exists $\varepsilon>0$ and $(a,b)\ni\theta_{0}$ such
that $R(\widetilde{\theta},\theta_{0})<\frac{1}{n}-\varepsilon$.
Consider $\pi_{m}=N(0,m)$. Then we have seen previously that the
Bayes risk is
\[
\int_{\mathbb{R}}R(\widehat{\theta}_{\pi_{m}},\theta)\pi_{m}(\theta)\,d\theta=\frac{1}{n+m^{-1}},
\]
and
\[
\frac{1}{n}-\int_{\mathbb{R}}R(\widehat{\theta}_{\pi_{m}},\theta)\pi_{m}(\theta)\,d\theta=\frac{1}{n}-\frac{1}{n+m^{-1}}\asymp\frac{1}{m}.
\]
Then
\begin{align*}
\frac{1}{n}-\int_{\mathbb{R}}R(\widetilde{\theta},\theta)\pi_{m}(\theta)\,d\theta & =\int_{\mathbb{R}}\left(\frac{1}{n}-R(\widetilde{\theta},\theta)\right)\pi_{m}(\theta)\,d\theta\\
 & =\int_{(a,b)}\left(\frac{1}{n}-R(\widetilde{\theta},\theta)\right)\pi_{m}(\theta)\,d\theta+\int_{\mathbb{R}\backslash(a,b)}\left(\frac{1}{n}-R(\widetilde{\theta},\theta)\right)\pi_{m}(\theta)\,d\theta\\
 & \geq\epsilon\int_{(a,b)}\pi_{m}(\theta)\,d\theta\\
 & =\epsilon\mathbb{P}\left(\frac{a}{\sqrt{m}}<N(0,1)<\frac{b}{\sqrt{m}}\right)\\
 & =\epsilon\int_{\frac{a}{\sqrt{m}}}^{\frac{b}{\sqrt{m}}}\frac{1}{\sqrt{2\pi}}e^{-x^{2}/2}\,dx\\
 & \asymp\frac{1}{\sqrt{m}}.
\end{align*}
So there exists $m$ sufficiently large such that
\[
\frac{1}{n}-\int_{\mathbb{R}}R(\widehat{\theta}_{\pi_{m}},\theta)\pi_{m}(\theta)\,d\theta<\frac{1}{n}-\int_{\mathbb{R}}R(\widetilde{\theta},\theta)\pi_{m}(\theta)\,d\theta,
\]
which is equivalent to 
\[
\int_{\mathbb{R}}R(\widetilde{\theta},\theta)\pi_{m}(\theta)\,d\theta<\int_{\mathbb{R}}R(\widehat{\theta}_{\pi_{m}},\theta)\pi_{m}(\theta)\,d\theta.
\]
This is a contradiction as $\widehat{\theta}_{\pi_{m}}$ achieves
the smallest Bayes risk by definition. 
\end{proof}
%

\subsection{James-Stein estimator}

Let $X_{1},\ldots,X_{n}\overset{\mathrm{i.i.d.}}{\sim}N(\theta,I_{p})$
where $\theta\in\mathbb{R}^{p}$, and consider the loss function $L(\widehat{\theta},\theta)=\Vert\widehat{\theta}-\theta\Vert^{2}$.
Admissibility of $\overline{X}$ in higher dimensions is much more
difficult. For $p=2$, admissibility of $\overline{X}$ was proved
by Stein. The proof involves constructing much more complex priors.
Surprisingly, $\overline{X}$ is no longer admissible if $p\geq3$. 
\begin{defn}
The James-Stein estimator is defined as
\[
\widehat{\theta}_{\text{JS}}=\left(1-\frac{p-2}{n\Vert\overline{X}\Vert^{2}}\right)\overline{X}.
\]
\end{defn}

It turns out that in $p\geq3$, the James-Stein estimator has a strictly
smaller risk than $\overline{X}$. Since $\overline{X}$ is minimax,
so is $\widehat{\theta}_{\text{JS}}$. This shows that minimax estimators
are not necessarily admissible, and also that minimax estimators are
not unique. 
\begin{lem}
[Stein's identity] If $Z\sim N(0,I_{p})$ then we have 
\[
\mathbb{E}\langle Z,g(Z)\rangle=\mathbb{E}\langle\nabla,g(Z)\rangle
\]
 for any $g\in C^{1}(\mathbb{R}^{p};\mathbb{R}^{p})$ for which both
expectations exist. 
\end{lem}

\begin{proof}
We will prove this in the $p=1$ case only. Let $\phi(x)=\frac{1}{\sqrt{2\pi}}e^{-x^{2}/2}$.
Then $\phi'(x)=-x\phi(x)$, so 
\[
\mathbb{E}[Zg(Z)]=\int xg(x)\phi(x)\,dx=-\int\phi'(x)g(x)\,dx=\int\phi(x)g'(x)\,dx=\mathbb{E}[g'(Z)]
\]
where the second last equality follows from integration by parts. 
\end{proof}
%
\begin{thm}
Let $p\geq3$. Then $\mathbb{E}_{\theta}\Vert\hat{\theta}_{\mathrm{JS}}-\theta\Vert^{2}<\frac{p}{n}=\mathbb{E}_{\theta}\Vert\overline{X}-\theta\Vert^{2}$
for all $\theta\in\mathbb{R}^{p}$.
\end{thm}

\begin{proof}
Let us write $\overline{X}=\frac{1}{\sqrt{n}}(\mu+Z)$ where $Z\sim N(0,I_{p})$
and $\mu=\sqrt{n}\theta$. We have
\begin{align*}
\mathbb{E}_{\theta}\Vert\hat{\theta}_{\mathrm{JS}}-\theta\Vert^{2} & =\mathbb{E}_{\theta}\left\Vert \left(1-\frac{p-2}{n\Vert\overline{X}\Vert^{2}}\right)\overline{X}-\theta\right\Vert ^{2}\\
 & =\mathbb{E}_{\theta}\left\Vert \overline{X}-\theta-\frac{p-2}{n\Vert\overline{X}\Vert^{2}}\overline{X}\right\Vert ^{2}\\
 & =\mathbb{E}_{\theta}\left\Vert \frac{1}{\sqrt{n}}Z-\frac{p-2}{\Vert\mu+Z\Vert^{2}}\frac{1}{\sqrt{n}}(\mu+Z)\right\Vert ^{2}\\
 & =\frac{1}{n}\mathbb{E}_{\theta}\left\Vert Z-\frac{p-2}{\Vert\mu+Z\Vert^{2}}(\mu+Z)\right\Vert ^{2}\\
 & =\frac{1}{n}\left(\mathbb{E}_{\theta}\Vert Z\Vert^{2}+\mathbb{E}_{\theta}\frac{(p-2)^{2}}{\Vert\mu+Z\Vert^{2}}-2(p-2)\mathbb{E}_{\theta}\left\langle Z,\frac{\mu+Z}{\Vert\mu+Z\Vert^{2}}\right\rangle \right).
\end{align*}
We have $\mathbb{E}_{\theta}\left\langle Z,\frac{\mu+Z}{\Vert\mu+Z\Vert^{2}}\right\rangle =\mathbb{E}_{\theta}\langle Z,g(Z)\rangle$,
where $g_{j}(z)=\frac{\mu_{j}+z_{j}}{\Vert\mu+z\Vert^{2}}$. Computing
the derivative gives
\[
\frac{\partial}{\partial z_{j}}g_{j}(z)=\frac{\Vert\mu+z\Vert^{2}-2(\mu_{j}+z_{j})^{2}}{\Vert\mu+z\Vert^{4}},
\]
so by Stein's identity we have
\begin{align*}
\mathbb{E}_{\theta}\left\langle Z,\frac{\mu+Z}{\Vert\mu+Z\Vert^{2}}\right\rangle  & =\sum_{j=1}^{p}\mathbb{E}_{\theta}\frac{\partial}{\partial z_{j}}g_{j}(Z)\\
 & =\sum_{j=1}^{p}\mathbb{E}_{\theta}\frac{\Vert\mu+Z\Vert^{2}-2(\mu_{j}+Z_{j})^{2}}{\Vert\mu+Z\Vert^{4}}\\
 & =\mathbb{E}_{\theta}\frac{(p-2)\Vert\mu+Z\Vert^{2}}{\Vert\mu+Z\Vert^{4}}\\
 & =\mathbb{E}_{\theta}\frac{p-2}{\Vert\mu+Z\Vert^{2}}.
\end{align*}
Then 
\begin{align*}
\mathbb{E}_{\theta}\Vert\hat{\theta}_{\mathrm{JS}}-\theta\Vert^{2} & =\frac{1}{n}\left(\mathbb{E}_{\theta}\Vert Z\Vert^{2}+\mathbb{E}_{\theta}\frac{(p-2)^{2}}{\Vert\mu+Z\Vert^{2}}-2(p-2)\mathbb{E}_{\theta}\left\langle Z,\frac{\mu+Z}{\Vert\mu+Z\Vert^{2}}\right\rangle \right).\\
 & =\frac{1}{n}\left(p-(p-2)^{2}\mathbb{E}_{\theta}\frac{1}{\Vert\mu+Z\Vert^{2}}\right)\\
 & <\frac{p}{n}.
\end{align*}
\end{proof}
%
We see that 
\[
R(\hat{\theta}_{\text{JS}},\theta)=\frac{1}{n}\left(p-(p-2)^{2}\mathbb{E}_{\theta}\frac{1}{\Vert\sqrt{n}\theta+Z\Vert^{2}}\right),\quad Z\sim N(0,1).
\]
By symmetry, $R(\hat{\theta}_{\text{JS}},\theta)$ depends only on
the magnitude of $\theta$, and we have 
\[
\lim_{\Vert\theta\Vert\to\infty}R(\hat{\theta}_{\text{JS}},\theta)=\frac{p}{n},\quad\text{so }\sup_{\theta\in\mathbb{R}^{p}}R(\hat{\theta}_{\text{JS}},\theta)=\sup_{\theta\in\mathbb{R}^{p}}R(\overline{X},\theta).
\]
Next, we consider the shrinkage estimator. Again, let $X_{1},\ldots,X_{n}\overset{\mathrm{iid}}{\sim}N(\theta,I_{p})$
where $\theta\in\mathbb{R}^{p}$ and consider the loss function $L(\hat{\theta},\theta)=\Vert\hat{\theta}-\theta\Vert^{2}$.
Let $\hat{\theta}_{c}=c\overline{X}$. We want to find $c^{\ast}$
which minimizes $R(\hat{\theta}_{c},\theta)$. We have
\[
R(\hat{\theta}_{c},\theta)=\mathbb{E}_{\theta}\Vert c\overline{X}-c\theta\Vert^{2}+\Vert c\theta-\theta\Vert^{2}=c^{2}\frac{p}{n}+(c-1)^{2}\Vert\theta\Vert^{2},
\]
so
\[
\frac{d}{dc}R(\hat{\theta}_{c},\theta)=2c\frac{p}{n}+2(c-1)\Vert\theta\Vert^{2}\implies c^{\ast}=\frac{\Vert\theta\Vert^{2}}{\frac{p}{n}+\Vert\theta\Vert^{2}}.
\]
The estimator 
\[
\hat{\theta}_{c^{\ast}}=\frac{\Vert\theta\Vert^{2}}{\frac{p}{n}+\Vert\theta\Vert^{2}}\overline{X}=\left(1-\frac{p}{p+n\Vert\theta\Vert^{2}}\right)\overline{X}
\]
is called the optimal shrinkage estimator. The risk function is
\[
R(\hat{\theta}_{c^{\ast}},\theta)=\left(\frac{b}{a+b}\right)^{2}a+\left(\frac{a}{a+b}\right)^{2}b=\frac{b^{2}a+a^{2}b}{(a+b)^{2}}=\frac{ab}{a+b}=\frac{\frac{p}{n}\Vert\theta\Vert^{2}}{\frac{p}{n}+\Vert\theta\Vert^{2}}\leq\min\left(\frac{p}{n},\Vert\theta\Vert^{2}\right),
\]
where $a=\frac{p}{n}$ and $b=\Vert\theta\Vert^{2}$. The optimal
shrinkage estimator is also called the oracle estimator because it
depends on the true parameter $\theta$ which is unknown. The optimal
shrinkage estimator is nearly optimal in the following sense.
\begin{thm}
[Oracle inequality] We have
\[
R(\hat{\theta}_{\mathrm{JS}},\theta)\leq\inf_{c}R(\hat{\theta}_{c},\theta)+\frac{2}{n}.
\]
\end{thm}

\begin{proof}
We have
\[
R(\hat{\theta}_{\mathrm{JS}},\theta)=\frac{1}{n}\left(p-(p-2)^{2}\mathbb{E}\frac{1}{\Vert\sqrt{n}\theta+Z\Vert^{2}}\right),\quad Z\sim N(0,1).
\]
Note that $Z+\sqrt{n}\theta\sim\chi_{p,n\Vert\theta\Vert^{2}}^{2}\overset{d}{=}\chi_{p+2N}^{2}$
where $N\sim\text{Poisson}\left(\frac{n\Vert\theta\Vert^{2}}{2}\right)$,
so
\[
\mathbb{E}\frac{1}{\Vert\sqrt{n}\theta+Z\Vert^{2}}=\mathbb{E}\frac{1}{\chi_{p+2N}^{2}}=\mathbb{E}\frac{1}{p+2N-2}\geq\frac{1}{\mathbb{E}(p+2N-2)}=\frac{1}{p+n\Vert\theta\Vert^{2}-2}.
\]
Therefore, 
\begin{align*}
R(\hat{\theta}_{\mathrm{JS}},\theta) & \leq\frac{1}{n}\left(p-\frac{(p-2)^{2}}{p+n\Vert\theta\Vert^{2}-2}\right)\\
 & =\frac{1}{n}\frac{np\Vert\theta\Vert^{2}+2(p-2)}{p+n\Vert\theta\Vert^{2}-2}\\
 & =\frac{1}{n}\left(2+\frac{(p-2)\Vert\theta\Vert^{2}}{p-2+n\Vert\theta\Vert^{2}}\right)\\
 & =\frac{2}{n}+\frac{\frac{p-2}{n}\Vert\theta\Vert^{2}}{\frac{p-2}{n}+\Vert\theta\Vert^{2}}\\
 & \leq\frac{2}{n}+R(\hat{\theta}_{c^{\ast}},\theta).
\end{align*}
\end{proof}
%
It is worth noting that $\frac{2}{n}$ is independent of the dimension
$p$, so in high dimensions the optimal shrinkage estimator performs
very closely to the James-Stein estimator. 

\section{Nonparametric Estimation}

\subsection{The problem of density estimation}

Let $(\varphi_{j})_{j=1}^{\infty}$ be a the orthonormal basis for
$L^{2}([0,1])$ given by $\{1,\sqrt{2}\sin(2\pi j\cdot),\sqrt{2}\cos(2\pi j\cdot):j=1,2,\ldots\}$.
Given $\alpha>0$, we define the Sobolev ball
\[
S_{\alpha}(R)=\left\{ f\in L^{2}([0,1]),\,f=\sum_{j}\theta_{j}\varphi_{j},\,f\geq0,\:\int f=1,\,\sum_{j}j^{2\alpha}\theta_{j}^{2}\leq R^{2}\right\} .
\]
The parameter $\alpha$ is called the smoothness parameter because
for a function $f\in C^{\alpha}[0,1]$ with Fourier series expansion
$f(x)=a_{0}+\sum_{j=1}^{\infty}\left(a_{j}\sin(2\pi jx)+b_{j}\cos(2\pi jx)\right)$,
we have
\[
\Vert f^{(\alpha)}\Vert^{2}=a_{0}^{2}+\frac{1}{2}\sum_{j=1}^{\infty}(2\pi j)^{2\alpha}(a_{j}^{2}+b_{j}^{2}).
\]
Let $f_{1},\ldots,f_{n}\overset{\text{iid}}{\sim}f\in S_{\alpha}(R)$
with loss function $\Vert\widehat{f}-f\Vert^{2}=\int|\widehat{f}(x)-f(x)|^{2}\,dx$.
Then estimating $f$ is equivalent to estimating the coefficients
$\theta_{j}$. Since $\theta_{j}=\int f\varphi_{j}$, by the central
limit theorem, we have
\[
\frac{1}{n}\sum_{i=1}^{n}\int f_{i}\varphi_{j}\overset{\text{approx}}{\sim}N\left(\theta_{j},\frac{\Var(\int f_{j}\varphi_{j})}{n}\right).
\]
However, instead of considering density estimation directly, we can
consider an asymptotically equivalent model known as the Gaussian
sequence model. Most nonparametric models such as the white noise
model and nonparametric regression are asymptotically equivalent to
the Gaussian sequence model. 

\subsection{Gaussian sequence model}

In the Gaussian sequence model, we have observations $X_{j}=\theta_{j}+\frac{1}{\sqrt{n}}Z_{j}$
where $Z_{j}\overset{\text{iid}}{\sim}N(0,1)$ for $j=1,2,\ldots$
and $\theta=(\theta_{j})\in\Theta_{\alpha}(R)=\{\theta:\sum_{j}j^{2\alpha}\theta_{j}^{2}<\infty\}$.
Consider the estimator
\[
\hat{\theta}_{j}=\begin{cases}
X_{j} & \text{if }j\leq k,\\
0 & \text{if }j>k.
\end{cases}
\]
Then the risk is
\begin{align*}
\mathbb{E}\Vert\hat{\theta}_{j}-\theta\Vert^{2} & =\sum_{j=1}^{k}\mathbb{E}(X_{j}-\theta_{j})^{2}+\sum_{j>k}\theta_{j}^{2}=\frac{k}{n}+\sum_{j>k}\theta_{j}^{2}\leq\frac{k}{n}+\frac{R^{2}}{k^{2\alpha}}.
\end{align*}
If $k$ is small then $R^{2}k^{-2\alpha}$ is large while if $k$
is large then $kn^{-1}$ is large, so the threshold $k$ should be
chosen to balance the two terms. Setting the terms to be equal, we
find that $k\asymp n^{\frac{1}{2\alpha+1}}$. With this choice, we
have
\[
\mathbb{E}\Vert\hat{\theta}_{j}-\theta\Vert^{2}\leq Cn^{-\frac{2\alpha}{2\alpha+1}}.
\]
This rate is optimal in the sense that the minimax risk is
\[
\inf_{\hat{\theta}}\sup_{\theta\in\Theta_{\alpha}(R)}\mathbb{E}_{\theta}\Vert\hat{\theta}-\theta\Vert^{2}=(1+o(1))C_{\alpha,R}n^{-\frac{2\alpha}{\alpha+1}},
\]
where $C_{\alpha,R}$ is called the Pinsker constant. However, an
issue with this estimator is that the choice of the threshold $k$
requires knowledge of $\alpha$ which is not known. 

Let us now consider the blockwise James-Stein estimator which does
not require knowledge of the smoothness parameter. Let $\{1,2,\ldots,n\}=B_{1}\cup\cdots\cup B_{m}$,
where $B_{\ell}=\{3^{\ell-1}+1,\ldots,3^{\ell}\}$ and $m=\log_{3}n$.
Consider the estimator
\[
\hat{\theta}=\left(\begin{array}{c}
\hat{\theta}_{B_{1}}\\
\hat{\theta}_{B_{2}}\\
\vdots\\
\hat{\theta}_{B_{m}}\\
0
\end{array}\right),\quad\hat{\theta}_{B_{\ell}}=\left(1-\frac{|B_{\ell}|-2}{n\Vert X_{\ell}\Vert^{2}}\right)X_{\ell}.
\]
Then by the oracle inequality we have
\begin{align*}
\mathbb{E}\Vert\hat{\theta}_{j}-\theta\Vert^{2} & =\sum_{j=1}^{m}\mathbb{E}\Vert\hat{\theta}_{B_{\ell}}-\theta_{B_{\ell}}\Vert^{2}+\sum_{j>m}\theta_{j}^{2}\leq\sum_{\ell=1}^{m}\min\left(\frac{|B_{\ell}|}{n},\Vert\theta_{B_{\ell}}\Vert^{2}\right)+\frac{2m}{n}+\sum_{j>n}\theta_{j}^{2}.
\end{align*}
Note that
\[
\frac{2m}{n}=\frac{\log_{3}n}{n}=o\left(n^{-\frac{2\alpha}{\alpha+1}}\right)\quad\text{and}\quad\sum_{j>n}\theta_{j}^{2}\leq\frac{R^{2}}{n^{2\alpha}}=o\left(n^{-\frac{2\alpha}{\alpha+1}}\right).
\]
Choosing $L=\min\{\ell\in\mathbb{N}:3^{\ell}\geq n^{\frac{1}{2\alpha+1}}\}$,
we have $3^{L-1}\leq n^{\frac{1}{2\alpha+1}}\leq3^{L}$, so
\begin{align*}
\sum_{\ell=1}^{m}\min\left(\frac{|B_{\ell}|}{n},\Vert\theta_{B_{\ell}}\Vert^{2}\right) & \leq\sum_{\ell=1}^{L}\frac{|B_{\ell}|}{n}+\sum_{\ell=L+1}^{m}\Vert\theta_{B_{\ell}}\Vert^{2}\\
 & \leq\sum_{\ell=1}^{L}\frac{2}{3n}3^{\ell}+\sum_{j>3^{L}}|\theta_{j}|^{2}\\
 & \apprle\frac{3^{L}}{n}+\sum_{j>n^{\frac{1}{2\alpha+1}}}|\theta_{j}|^{2}\\
 & \apprle\frac{n^{\frac{1}{2\alpha+1}}}{n}+\left(n^{\frac{1}{2\alpha+1}}\right)^{-2\alpha}\\
 & \apprle n^{-\frac{2\alpha}{2\alpha+1}}.
\end{align*}
Adding up all the terms, we get $\mathbb{E}\Vert\hat{\theta}_{j}-\theta\Vert^{2}\apprle n^{-\frac{2\alpha}{2\alpha+1}}.$
So the block-wise James-Stein estimator also achieves the minimax
rate. 

\section{Hypothesis Testing}

\subsection{Optimal testing error for a simple two-point test}

In a two-point test, the null hypothesis is $H_{0}:X\sim P$ and the
alternative hypothesis is $H_{1}:X\sim Q$. The testing function is
a function $\phi:X\mapsto\{0,1\}$. We define the Type-I error to
be $P\phi:=\mathbb{E}_{X\sim P}\phi(X)$ and the Type-II error to
be $Q(1-\phi):=\mathbb{E}_{X\sim Q}(1-\phi(X))$. The testing error
is then $P\phi+Q(1-\phi)$, and the optimal testing error is $\inf_{\phi}(P\phi+Q(1-\phi))$. 
\begin{defn}
The total variation distance between two probability measures $P$
and $Q$ is
\[
\mathsf{TV}(P,Q)=\sup_{B}\left|P(B)-Q(B)\right|
\]
where the supremum is taken over all measurable sets $B$. 
\end{defn}

\begin{thm}
Suppose that $P$ and $Q$ are absolutely continuous with respect
to some dominating measure with density function $p$ and $q$. Then
\begin{align*}
\mathsf{TV}(P,Q) & =P(\{p(x)>q(x)\})-Q(\{p(x)>q(x)\})=\frac{1}{2}\int|p-q|=1-\int p\land q.
\end{align*}
\end{thm}

\begin{proof}
Let $A:=\{p(x)>q(x)\}$. Then $\mathsf{TV}(P,Q)\geq P(A)-Q(A)$ by
definition. For all measurable sets $B$, we have
\begin{align*}
|P(B)-Q(B)| & =\left|\int_{B}(p-q)\right|\\
 & =\left|\int_{B\cap A}(p-q)+\int_{B\cap A^{c}}(p-q)\right|\\
 & =\left|\int_{B\cap A}(p-q)-\int_{B\cap A^{c}}(q-p)\right|\\
 & \leq\max\left(\int_{B\cap A}(p-q),\int_{B\cap A^{c}}(q-p)\right)\\
 & \leq\max\left(\int_{A}(p-q),\int_{A^{c}}(q-p)\right)\\
 & =\max\left(P(A)-Q(A),Q(A^{c})-P(A^{c})\right)\\
 & =P(A)-Q(A).
\end{align*}
Taking supremum over $B$ proves the first equality. For the second,
we note that
\begin{align*}
\frac{1}{2}\int|p-q| & =\frac{1}{2}\int_{A}(p-q)+\frac{1}{2}\int_{A^{c}}(q-p)\\
 & =\frac{1}{2}(P(A)-Q(A)+Q(A^{c})-P(A^{c}))\\
 & =P(A)-Q(A).
\end{align*}
For the third equality, we note that
\[
\int p\land q=\int_{A}q+\int_{A^{c}}p=Q(A)+1-P(A)=1-(P(A)-Q(A)),
\]
and rearranging gives the desired equality. 
\end{proof}
%
\begin{thm}
[Neyman-Pearson lemma] The optimal testing error is given by
\[
\inf_{\phi}(P\phi+Q(1-\phi))=\int p\land q=1-\mathsf{TV}(P,Q)
\]
which is achieved by the likelihood ratio test $\phi^{\ast}(X)=\mathbb{{1}}_{\{p\leq q\}}(X)$. 
\end{thm}

\begin{proof}
For all $\phi$, we have $P\phi+Q(1-\phi)=\int p\phi+\int q(1-\phi)=\int p\phi+q(1-\phi)\geq\int\min(p,q)$,
so the optimal testing error is bounded below by $1-\mathsf{TV}(P,Q)=\int\min(p,q)$.
Conversely, we have
\[
\inf_{\phi}(P\phi+Q(1-\phi))\leq P\phi^{\ast}+Q(1-\phi^{\ast})=P(A^{c})+Q(A)=1-\mathsf{TV}(P,Q)=\int\min(p,q),
\]
where $A:=\{p(x)>q(x)\}$. 
\end{proof}
%

\subsection{Le Cam's two point method}

Le Cam's two point method gives us a way to lower bound the minimax
risk. The idea is to fix two points $\theta_{0},\theta_{1}\in\Theta$
and put the prior $\pi=\frac{1}{2}\delta_{\theta_{0}}+\frac{1}{2}\delta_{\theta_{1}}$
on $\theta$, so that
\begin{align*}
\sup_{\theta\in\Theta}\mathbb{E}(\widehat{\theta}-\theta)^{2} & \geq\frac{1}{2}\mathbb{E}_{\theta_{0}}(\widehat{\theta}-\theta_{0})^{2}+\frac{1}{2}\mathbb{E}_{\theta_{1}}(\widehat{\theta}-\theta_{1})^{2}\\
 & =\frac{1}{2}\int\left[(\widehat{\theta}-\theta_{0})^{2}p_{\theta_{0}}+(\widehat{\theta}-\theta_{1})^{2}p_{\theta_{1}}\right]\\
 & \geq\frac{1}{2}\int\left[(\widehat{\theta}-\theta_{0})^{2}+(\widehat{\theta}-\theta_{1})^{2}\right]p_{\theta_{0}}\land p_{\theta_{1}}\\
 & \geq\frac{(\theta_{0}-\theta_{1})^{2}}{4}\int p_{\theta_{0}}\land p_{\theta_{1}},
\end{align*}
where the last inequality follows from $(a+b)^{2}\leq2a^{2}+2b^{2}$
which holds for all $a,b\in\mathbb{R}$. Taking infimum gives the
following. 
\begin{thm}
[Le Cam's two point method] If $L(\widehat{\theta},\theta)=(\widehat{\theta}-\theta)^{2}$,
then for all $\theta_{0},\theta_{1}\in\Theta$, we have
\[
\inf_{\widehat{\theta}}\sup_{\theta\in\Theta}\mathbb{E}(\hat{\theta}-\theta)^{2}\geq\frac{(\theta_{0}-\theta_{1})^{2}}{4}\int\min(p_{\theta_{0}},p_{\theta_{1}}).
\]
\end{thm}

\begin{example}
Let $X_{1},\ldots,X_{n}\overset{\text{iid}}{\sim}N(\theta,1)$ where
$\theta\in\mathbb{R}$. We already know from before that the minimax
risk is exactly $\frac{1}{n}$. We will re-derive this lower bound
(up to a constant). Let $\theta_{0}=0$ and $\theta_{1}=\frac{1}{\sqrt{n}}$.
Then by Le Cam's lemma, 
\begin{align*}
\inf_{\widehat{\theta}}\sup_{\theta\in\Theta}\mathbb{E}(\hat{\theta}-\theta)^{2} & \geq\frac{1}{4n}\int p_{0}^{n}\land p_{1/\sqrt{n}}^{n}\\
 & =\frac{1}{4n}\inf_{\phi}\left(P_{0}^{n}\phi+P_{\frac{1}{\sqrt{n}}}^{n}(1-\phi)\right)\\
 & \geq\frac{1}{4n}P_{0}^{n}(\{p_{0}^{n}\leq p_{\frac{1}{\sqrt{n}}}^{n}\})\\
 & =\frac{1}{4n}P_{0}^{n}\left(\left\{ \prod_{i=1}^{n}\frac{e^{-\frac{1}{2}(X_{i}-\frac{1}{\sqrt{n}})^{2}}}{e^{-\frac{1}{2}X_{i}^{2}}}\geq1\right\} \right)\\
 & =\frac{1}{4n}P_{0}^{n}\left(\left\{ \sum_{i=1}^{n}\left[\left(X_{i}-\frac{1}{\sqrt{n}}\right)^{2}-X_{i}^{2}\right]<0\right\} \right)\\
 & =\frac{1}{4n}P_{0}^{n}\left(\left\{ \frac{1}{\sqrt{n}}\sum_{i=1}^{n}X_{i}>\frac{1}{2}\right\} \right)\\
 & =\frac{1}{4n}\mathbb{P}(N(0,1)>\frac{1}{2}),
\end{align*}
so we have the lower bound
\[
\inf_{\hat{\theta}}\sup_{\theta\in\mathbb{R}}\mathbb{E}_{\theta}(\widehat{\theta}-\theta)^{2}\geq Cn^{-1}
\]
where $C:=\mathbb{P}(N(0,1)>\frac{1}{2})/4$. 
\end{example}

\begin{example}
Let $X_{1},\ldots,X_{n}\overset{\text{iid}}{\sim}\text{Unif}(0,\theta)$
where $\theta\in[0,1]$. Let $\widehat{\theta}=X_{(n)}$. Then
\begin{align*}
\mathbb{E}_{\theta}(\widehat{\theta}-\theta)^{2} & =\int_{0}^{\theta}(t-\theta)^{2}\theta^{-n}nt^{n-1}\,dt=\frac{2\theta^{2}}{(n+1)(n+2)},
\end{align*}
so 
\[
\sup_{\theta\in[0,1]}\mathbb{E}_{\theta}(\widehat{\theta}-\theta)^{2}=\frac{2}{(n+1)(n+2)}=O(n^{-2}).
\]
This gives an upper bound on the minimax rate. For the lower bound,
we use Le Cam's two point method. Let $\theta_{1}=1$ and $\theta_{2}=1-\frac{1}{n}$.
Then
\begin{align*}
\inf_{\widehat{\theta}}\sup_{\theta\in[0,1]}\mathbb{E}_{\theta}(\hat{\theta}-\theta)^{2} & \geq\frac{1}{4n^{2}}\int p_{1}^{n}\land p_{1-\frac{1}{n}}^{n}\\
 & \geq\frac{1}{8n^{2}}\left(\int\sqrt{p_{1}^{n}p_{1-\frac{1}{n}}^{n}}\right)^{2}\\
 & =\frac{1}{8n^{2}}\left(\int\sqrt{p_{1}p_{1-\frac{1}{n}}}\right)^{2n}\\
 & =\frac{1}{8n^{2}}\left(\int_{0}^{1-\frac{1}{n}}\sqrt{\frac{1}{1-\frac{1}{n}}}\right)^{2n}\\
 & =\frac{1}{8n^{2}}\left(1-\frac{1}{n}\right)^{n},
\end{align*}
where for the second inequality we used $\int p\land q\geq\frac{1}{2}\left(\int\sqrt{pq}\right)^{2}$.
So we have
\[
\inf_{\hat{\theta}}\sup_{\theta\in\mathbb{R}}\mathbb{E}_{\theta}(\hat{\theta}-\theta)^{2}\asymp n^{-2}.
\]
\end{example}


\subsection{Lower bound for Gaussian sequence model}

Now we want to prove the minimax lower bound for the Gaussian sequence
problem. Recall that we have observations $X_{j}\sim N(\theta_{j},\frac{1}{n})$,
$j\in\mathbb{N}$, where $\theta\in\Theta_{\alpha}(R)=\{\theta:\sum_{j=1}^{\infty}j^{2\alpha}\theta_{j}^{2}\leq R^{2}\}$.
We want to show that
\[
\inf_{\hat{\theta}}\sup_{\theta\in\Theta_{\alpha}(R)}\mathbb{E}_{\theta}\Vert\hat{\theta}-\theta\Vert^{2}\geq cn^{\frac{-2\alpha}{2\alpha+1}}.
\]
Let us define $\Theta_{0}=\{\theta:\theta_{j}\in\{0,\frac{1}{\sqrt{n}}\}\text{ for }j=1,\ldots,k,\theta_{j}=0\text{ for }j>k\}$.
Then $|\Theta_{0}|=2^{k}$. To ensure that $\Theta_{0}\subseteq\Theta_{\alpha}(R)$,
we note that
\[
\sum_{j=1}^{\infty}j^{2\alpha}\theta_{j}^{2}\leq\frac{1}{n}\sum_{j=1}^{k}j^{2\alpha}\leq\frac{k^{2\alpha+1}}{n}.
\]
Solving $\frac{k^{2\alpha+1}}{n}\leq R^{2}$ gives $k\asymp n^{\frac{1}{2\alpha+1}}$.
Then
\begin{align*}
\inf_{\widehat{\theta}}\sup_{\theta\in\Theta_{\alpha}(R)}\mathbb{E}_{\theta}\Vert\widehat{\theta}-\theta\Vert^{2} & \geq\inf_{\widehat{\theta}}\sup_{\theta\in\Theta_{0}}\mathbb{E}_{\theta}\Vert\widehat{\theta}-\theta\Vert^{2}\\
 & \geq\inf_{\widehat{\theta}}\operatornamewithlimits{ave}_{\theta\in\Theta_{0}}\mathbb{E}_{\theta}\Vert\widehat{\theta}-\theta\Vert^{2}\\
 & \geq\inf_{\widehat{\theta}}\operatornamewithlimits{ave}_{\theta\in\Theta_{0}}\sum_{j=1}^{k}(\widehat{\theta}_{j}-\theta_{j})^{2}\\
 & =\inf_{\widehat{\theta}}\sum_{j=1}^{k}\operatornamewithlimits{ave}_{\theta\in\Theta_{0}}\mathbb{E}_{\theta}(\widehat{\theta}_{j}-\theta_{j})^{2}\\
 & =\inf_{\widehat{\theta}}\sum_{j=1}^{k}\operatornamewithlimits{ave}_{\theta_{-j}}\left[\frac{1}{2}\mathbb{E}_{(\theta_{-j},\theta_{j}=0)}(\widehat{\theta}_{j})^{2}+\frac{1}{2}\mathbb{E}_{(\theta_{-j},\theta_{j}=\frac{1}{\sqrt{n}})}(\widehat{\theta}_{j}-\frac{1}{\sqrt{n}})^{2}\right]\\
 & \geq\sum_{j=1}^{k}\operatornamewithlimits{ave}_{\theta_{-j}}\inf_{\widehat{\theta}_{j}}\left[\frac{1}{2}\mathbb{E}_{(\theta_{-j},\theta_{j}=0)}(\widehat{\theta}_{j})^{2}+\frac{1}{2}\mathbb{E}_{(\theta_{-j},\theta_{j}=\frac{1}{\sqrt{n}})}(\widehat{\theta}_{j}-\frac{1}{\sqrt{n}})^{2}\right].
\end{align*}
Using the same argument as in the proof of Le Cam's method to estimate
the infimum, we have
\[
\inf_{\widehat{\theta}}\sup_{\theta\in\Theta_{\alpha}(R)}\mathbb{E}_{\theta}\Vert\widehat{\theta}-\theta\Vert^{2}\geq\sum_{j=1}^{k}\operatornamewithlimits{ave}_{\theta_{-j}}\frac{1}{4n}\int p_{(\theta_{-j},\theta_{j}=0)}\land p_{(\theta_{-j},\theta_{j}=\frac{1}{\sqrt{n}})}.
\]
Since only the $\theta_{j}$ differs between the two hypotheses, this
is equivalent to testing $H_{0}:X_{j}\sim N(0,\frac{1}{n})$ versus
$H_{1}:X_{j}\sim N(\frac{1}{\sqrt{n}},\frac{1}{n})$. Therefore, letting
$P=N(0,\frac{1}{n})$ and $Q=N(\frac{1}{\sqrt{n}},\frac{1}{n})$,
we have
\[
\inf_{\widehat{\theta}}\sup_{\theta\in\Theta_{\alpha}(R)}\mathbb{E}_{\theta}\Vert\widehat{\theta}-\theta\Vert^{2}\geq\sum_{j=1}^{k}\operatornamewithlimits{ave}_{\theta_{-j}}\inf_{\phi}(P\phi+Q(1-\phi)).
\]
From an example above, we know that the optimal testing error for
this test satisfies 
\[
\inf_{\phi}(P\phi+Q(1-\phi))\geq c
\]
where $c=\mathbb{P}(N(0,1)>\frac{1}{2})$. Therefore, 
\[
\inf_{\widehat{\theta}}\sup_{\theta\in\Theta_{\alpha}(R)}\mathbb{E}_{\theta}\Vert\widehat{\theta}-\theta\Vert^{2}\geq\frac{ck}{4n}\asymp n^{-\frac{2\alpha}{2\alpha+1}}.
\]


\section{Uniformly Minimum-Variance Unbiased Estimators}

\subsection{UMVUEs}

We say that an estimator $\delta(X)$ is a uniformly minimum-variance
unbiased estimator (UMVUE) for $g(\theta)$ if $\mathbb{E}_{\theta}\delta(X)=g(\theta)$
for all $\theta\in\Theta$, and for any other estimator $\widetilde{\delta}(X)$
which is unbiased for $g(\theta)$, it holds that $\operatornamewithlimits{Var}_{\theta}(\delta(X))\leq\operatornamewithlimits{Var}_{\theta}(\widetilde{\delta}(X))$
for all $\theta\in\Theta$. 
\begin{thm}
[Lehmann-Scheffe] Suppose that $T$ is unbiased, sufficient, and
complete for $\theta$. If $\delta(X)=h(T(X))$ is unbiased for $g(\theta)$,
then
\begin{enumerate}
\item $\delta(X)$ is the only function of $T(X)$ which is unbiased for
$g(\theta)$,
\item $\delta(X)$ is the unique UMVUE for $g(\theta)$.
\end{enumerate}
\end{thm}

\begin{proof}
Suppose that $\widetilde{h}(T(X))$ is also unbiased for $g(\theta)$.
Then $\mathbb{E}_{\theta}(\widetilde{h}(T(X))-h(T(X)))$, so by completeness,
it follows that $\widetilde{h}(T(X))=h(T(X))$ almost surely. To show
the second statement, suppose that $\widetilde{\delta}(X)$ is unbiased
for $g(\theta)$. Then $\mathbb{E}(\widetilde{\delta}(X)|T(X))$ is
also unbiased, so $\mathbb{E}(\widetilde{\delta}(X)|T(X))=h(T(X))$.
By the Rao-Blackwell theorem, 
\[
\mathbb{E}_{\theta}\left(\mathbb{E}_{\theta}(\widetilde{\delta}(X)|T(X))-g(\theta)\right)^{2}\leq\mathbb{E}_{\theta}\left(\widetilde{\delta}(X)-g(\theta)\right)^{2}.
\]
The left hand side is $\operatornamewithlimits{Var}_{\theta}(h(T(X))$
and the right hand side is $\operatornamewithlimits{Var}_{\theta}(\widetilde{\delta}(X))$. 
\end{proof}
%
There are two methods to find UMVUEs:
\begin{enumerate}
\item Find $h(T(X))$ for some sufficient and complete $T(X)$ such that
$\mathbb{E}_{\theta}h(T(X))=g(\theta)$ for all $\theta\in\Theta$.
\item Find $\delta(X)$ such that $\mathbb{E}_{\theta}\delta(X)=g(\theta)$
for all $\theta\in\Theta$. The UMVUE is then $\mathbb{E}(\delta(X)|T(X))$
where $T(X)$ is sufficient and complete for $\theta$. 
\end{enumerate}
%
\begin{example}
Let $X_{1},\ldots,X_{n}\overset{\text{i.i.d.}}{\sim}\text{Ber}(\theta)$
where $\theta\in(0,1)$. 
\begin{enumerate}
\item UMVUE for $\theta$. The estimator $\widehat{\theta}=\overline{X}$
is unbiased, sufficient, and complete, it is the unique UMVUE.
\item UMVUE for $\theta^{2}$. Since $\mathbb{E}(X_{1}X_{2})=\theta^{2}$,
the estimator $\delta(X)=\mathbb{E}(X_{1}X_{2}|\sum_{i=1}^{n}X_{i})$
is unbiased. Since $T(X)$ is sufficient and complete for $\theta$,
it follows that $\delta(X)$ is the unique UMVUE for $\theta^{2}$. 
\end{enumerate}
\end{example}

\begin{example}
Let $X_{1},\ldots,X_{n}\overset{\text{i.i.d.}}{\sim}\text{Unif}(0,\theta)$
where $\theta>0$. To find the UMVUE for $\theta$, we start with
$\mathbb{E}_{\theta}(2X_{1})=\theta$, so the UMVUE is $\mathbb{E}_{\theta}(2X_{1}|X_{(n)})$.
Since $X_{1}|X_{(n)}\sim\frac{1}{n}\delta_{X_{(n)}}+\left(1-\frac{1}{n}\right)\text{Unif}(0,X_{(n)})$,
it follows that 
\begin{align*}
\mathbb{E}_{\theta}(2X_{1}|X_{(n)}) & =2\left(\frac{1}{n}X_{n}+\left(1-\frac{1}{n}\right)\frac{X_{n}}{2}\right)=\left(1+\frac{1}{n}\right)X_{(n)}.
\end{align*}
\end{example}

\begin{example}
Let $X_{1},\ldots,X_{n}\overset{\text{i.i.d.}}{\sim}N(\mu,\sigma^{2})$.
A complete sufficient statistic is $\left(\sum_{i=1}^{n}X_{i},\sum_{i=1}^{n}X_{i}^{2}\right)$,
which is equivalent to $\left(\overline{X},\sum_{i=1}^{n}(X_{i}-\overline{X})^{2}\right)$.
The UMVUE for $\mu$ is $\overline{X}$, and the UMVUE for $\sigma^{2}$
is $\frac{1}{n-1}\sum_{i=1}^{n}(X_{i}-\overline{X})^{2}$. 
\end{example}


\subsection{Stein's unbiased risk estimate}

Consider $X\sim N(\theta,I_{p})$ and let $\widehat{\theta}$ be an
estimator. We want to estimate the risk function $g(\theta)=\mathbb{E}_{\theta}\Vert\widehat{\theta}-\theta\Vert^{2}$.
Note that $X$ is complete and sufficient, so we want to write $\mathbb{E}_{\theta}\Vert\widehat{\theta}-\theta\Vert^{2}$
as the expectation of some function of $X$ which does not depend
on $\theta$. We have
\begin{align*}
\mathbb{E}_{\theta}\Vert\widehat{\theta}-\theta\Vert^{2} & =\mathbb{E}_{\theta}\Vert\widehat{\theta}-X+X-\theta\Vert^{2}\\
 & =\mathbb{E}_{\theta}\left(\Vert\widehat{\theta}-X\Vert^{2}+\Vert X-\theta\Vert^{2}+2\langle X-\theta,\widehat{\theta}-X\rangle\right)\\
 & =\mathbb{E}_{\theta}\left(\Vert\widehat{\theta}-X\Vert^{2}+\Vert X-\theta\Vert^{2}-2\langle X-\theta,X\rangle+2\langle X-\theta,\widehat{\theta}\rangle\right).
\end{align*}
Writing $X=\theta+Z$ where $Z\sim N(0,I_{p})$, we have $\mathbb{E}_{\theta}\Vert X-\theta\Vert^{2}=p,$
$\mathbb{E}_{\theta}\langle X-\theta,X\rangle=\mathbb{E}_{\theta}\langle Z,\theta+Z\rangle=p,$
and by Stein's lemma 
\begin{align*}
\mathbb{E}_{\theta}\langle X-\theta,\widehat{\theta}\rangle & =\mathbb{E}_{\theta}\langle Z,\widehat{\theta}(\theta+Z)\rangle=\mathbb{E}_{\theta}\sum_{j=1}^{p}\frac{\partial}{\partial z_{j}}\widehat{\theta}_{j}(\theta+Z)=\mathbb{E}_{\theta}\sum_{j=1}^{p}\frac{\partial}{\partial x_{j}}\widehat{\theta}_{j}.
\end{align*}
Therefore, $\mathbb{E}_{\theta}\Vert\widehat{\theta}-\theta\Vert^{2}=\mathbb{E}_{\theta}\operatorname{SURE}(\widehat{\theta})$
where 
\[
\operatorname{SURE}(\widehat{\theta})=\Vert\widehat{\theta}-X\Vert^{2}-p+2\sum_{j=1}^{p}\frac{\partial}{\partial x_{j}}\widehat{\theta}_{j}
\]
is called Stein's unbiased risk estimator. This is a UMVUE of the
risk function. 

\subsubsection{Applications}
\begin{enumerate}
\item Let $\widehat{\theta}_{c}=cX$. Then $\widehat{\theta}_{c^{\ast}}$
where $c^{\ast}=\argmin_{c}\operatorname{SURE}(cX)$ will give something
similar to the James-Stein estimator. 
\item Consider the linear model $Y\sim N(\theta,I_{n})$. Let $X\in\mathbb{R}^{n\times p}$
be a fixed full-rank design matrix and let $\widehat{\beta}=\argmin_{\beta}\Vert Y-X\beta\Vert^{2}$.
Then $\widehat{\theta}=X\widehat{\beta}=HY$ where $H=X(X^{T}X)^{-1}X^{T}$.
Then
\[
\sum_{i=1}^{n}\frac{\partial}{\partial y_{i}}\widehat{\theta}_{i}=\sum_{i=1}^{n}H_{ii}=\operatorname{Tr}(H)=p,
\]
which gives the Akaike information criterion
\[
\operatorname{SURE}(X\widehat{\beta})=\Vert Y-X\widehat{\beta}\Vert^{2}+2p-n.
\]
\end{enumerate}
%

\section{Equivariant Estimators}

\subsection{Location equivariance}

A family of distributions $(P_{\theta}:\theta\in\mathbb{R}^{p})$
is called a location model if $X\sim P_{\theta}\iff X-\theta\sim P_{0}$.
Assuming density exists, it is easy to verify that
\[
P_{\theta}(x)=f(x-\theta)
\]
where $f(x)=P_{0}(x)$. We say that an estimator $\widehat{\theta}$
is location equivariant if
\[
\widehat{\theta}(X_{1}+c,\ldots,X_{n}+c)=\widehat{\theta}(X_{1},\ldots,X_{n})+c.
\]
We assume that the loss function $L$ is location invariant, meaning
that 
\[
L(\widehat{\theta}+c,\theta+c)=L(\widehat{\theta},\theta).
\]
As a consequence of location invariance, we have $L(\widehat{\theta},\theta)=L(\widehat{\theta}-\theta,0)=:\rho(\widehat{\theta}-\theta)$,
so the loss function can be written as a function of the difference
between $\widehat{\theta}$ and $\theta$. 
\begin{thm}
For a location model $(P_{\theta}:\theta\in\mathbb{R}^{p})$, a location
invariant loss $L(\widehat{\theta},\theta)$, and a location equivariant
estimator $\widehat{\theta}$, the risk $\mathbb{E}_{\theta}L(\widehat{\theta},\theta)$
must be constant. 
\end{thm}

\begin{proof}
For all $\theta\in\mathbb{R}^{p}$, we have $\mathbb{E}_{\theta}L(\widehat{\theta},\theta)=\mathbb{E}_{\theta}\rho(\widehat{\theta}-\theta)=\mathbb{E}_{\theta}\rho(\widehat{\theta}(X_{1},\ldots,X_{n})-\theta)=\mathbb{E}_{\theta}\rho(\widehat{\theta}(X_{1}-\theta,\ldots,X_{n}-\theta))=\mathbb{E}_{0}\rho(\widehat{\theta}(X_{1},\ldots,X_{n}))$. 
\end{proof}
%

\subsubsection{Minimal risk equivariant estimator}

A minimal risk equivariant estimator (MREE) is an equivariant estimator
which has the lowest risk among all equivariant estimators. The strategy
to finding the MREE is to start with a pilot estimator $\delta_{0}(X)$
which can be any equivariant estimator. Then we set
\[
\delta(X)=\delta_{0}(X)-u(X).
\]
Note that $u(X+c)=\delta_{0}(X+c)-\delta(X+c)=\delta_{0}(X)+c-\delta(X)-c=\delta_{0}(X)-\delta(X)=u(X)$,
so $u(X_{1}+c,\ldots,X_{n}+c)=u(X_{1},\ldots,X_{n}).$ Taking $c=-X_{n}$,
we have
\[
u(X_{1},\ldots,X_{n})=u(X_{1}-X_{n},\ldots,X_{n-1}-X_{n},0).
\]
Let $Y=(X_{1}-X_{n},\ldots,X_{n-1}-X_{n})$ and write $u(X)=v(Y)$.
Then
\[
\mathbb{E}_{\theta}\rho(\delta(X)-\theta)=\mathbb{E}_{0}\rho(\delta(X))=\mathbb{E}_{0}\rho(\delta_{0}(X)-v(Y))=\mathbb{E}_{0}[\mathbb{E}_{0}(\rho(\delta_{0}(X)-v(Y))|Y)].
\]
Therefore, to minimize the risk we should choose
\[
v(Y)=\argmin_{a}\mathbb{E}_{0}(\rho(\delta_{0}(X)-a)|Y).
\]
In particular, when $\rho(t)=t^{2}$, $v(Y)=\mathbb{E}_{0}(\delta_{0}(X)|Y)$. 
\begin{example}
Suppose that $X_{1},\ldots,X_{n}\overset{\text{i.i.d.}}{\sim}N(\theta,1)$.
Let $\delta_{0}(X)=\overline{X}$. Then $\mathbb{E}_{0}(\overline{X}|Y)=\mathbb{E}_{0}(\overline{X})=0$,
so $\delta(X)=\overline{X}$ is the MREE. 
\end{example}

\begin{example}
Suppose that $X_{1},\ldots,X_{n}\overset{\text{i.i.d.}}{\sim}E(\theta,b)$.
The density is given by $p(x)=b^{-1}e^{-\frac{x-\theta}{b}}1_{\{x>\theta\}}$.
Let $\delta_{0}(X)=X_{(1)}$. Then since $Y$ is ancillary, $\mathbb{E}_{0}(X_{(1)}|Y)=\mathbb{E}_{0}(X_{(1)})$.
To evaluate this expectation, write $X_{i}=\theta+bZ_{i}$ where $Z_{i}$
has density $e^{-x}1_{\{x>0\}}$. Then $X_{(1)}=\theta+bZ_{(1)}$.
We have 
\[
\mathbb{P}(Z_{(1)}\leq t)=1-\mathbb{P}(Z_{(1)}>t)=1-\prod_{i=1}^{t}\mathbb{P}(Z_{i}>t)=1-e^{-nt}.
\]
The density of $Z_{(1)}$ is then $ne^{-nt}1_{\{t>0\}}$ which has
mean $\frac{1}{n}$. Therefore, $\mathbb{E}_{0}(X_{(1)})=\frac{b}{n}$,
which gives the MREE
\[
\delta(X)=X_{(1)}-\frac{b}{n}.
\]
\end{example}


\subsubsection{Pitman estimator}

Suppose that $\rho(t)=t^{2}$ and we have the pilot estimator $\delta_{0}(X)=X_{n}$.
Then $\delta(X)=X_{n}-\mathbb{E}_{0}(X_{n}|Y)$. Let 
\[
Z=\left(\begin{array}{c}
Y\\
X_{n}
\end{array}\right)=\left(\begin{array}{c}
X_{1}-X_{n}\\
\vdots\\
X_{n-1}-X_{n}\\
X_{n}
\end{array}\right).
\]
Then $X_{1}=Z_{1}+Z_{n},\ldots,X_{n-1}=Z_{n-1}+Z_{n}$, $X_{n}=Z_{n}$,
so 
\[
\frac{\partial X}{\partial Z}=\left(\begin{array}{cc}
I_{n-1} & 1_{n-1}^{T}\\
0 & 1
\end{array}\right),
\]
and writing $f(x)$ for $P_{0}(x)$, we have
\begin{align*}
P_{Z,0}(z_{1},\ldots,z_{n}) & =P_{X,0}(x_{1},\ldots,x_{n})\left|\frac{\partial X}{\partial Z}\right|\\
 & =P_{X,0}(x_{1},\ldots,x_{n})\\
 & =\prod_{i=1}^{n}f(x_{i})\\
 & =f(y_{1}+x_{n})\cdots f(y_{n-1}+x_{n})f(x_{n}).
\end{align*}
Then
\[
\mathbb{E}_{0}(X_{n}|Y)=\frac{\int tf(Y_{1}+t)\cdots f(Y_{n-1}+t)f(t)\,dt}{\int f(Y_{1}+t)\cdots f(Y_{n-1}+t)f(t)\,dt}.
\]
Using $X_{n}-\mathbb{E}_{0}(X_{n}|Y)$, we have
\begin{align*}
\delta(X) & =\frac{\int(X_{n}-t)f(Y_{1}+t)\cdots f(Y_{n-1}+t)f(t)\,dt}{\int f(Y_{1}+t)\cdots f(Y_{n-1}+t)f(t)\,dt}\\
 & =\frac{\int uf(X_{1}-u)\cdots f(X_{n}-u)\,du}{\int f(X_{1}-u)\cdots f(X_{n}-u)\,du}.
\end{align*}
The formula
\[
\delta(X)=\frac{\int u\prod_{i=1}^{n}f(X_{i}-u)\,du}{\int\prod_{i=1}^{n}f(X_{i}-u)\,du}
\]
 is called the Pitman estimator. 
\begin{example}
Let $X_{1},\ldots,X_{n}\overset{\text{i.i.d.}}{\sim}\text{Unif}(\theta-\frac{b}{2},\theta+\frac{b}{2})$.
Then the density is $P_{\theta}(x)=b^{-1}1_{\{\theta-\frac{b}{2}<x<\theta+\frac{b}{2}\}}.$
Then
\begin{align*}
\prod_{i=1}^{n}f(X_{i}-u) & =\prod_{i=1}^{n}b^{-1}1_{\{u-\frac{b}{2}<X_{i}<u+\frac{b}{2}\}}\\
 & =b^{-n}1_{\{u-\frac{b}{2}<X_{(1)}\}}1_{\{X_{(n)}<u+\frac{b}{2}\}}\\
 & =b^{-n}1_{\{X_{(n)}-\frac{b}{2}<u<X_{(1)}+\frac{b}{2}\}}.
\end{align*}
Therefore, the MREE is
\begin{align*}
\delta(X) & =\frac{\int_{X_{(n)}-\frac{b}{2}}^{X_{(1)}+\frac{b}{2}}u\,du}{\int_{X_{(n)}-\frac{b}{2}}^{X_{(1)}+\frac{b}{2}}\,du}=\frac{1}{2}\frac{(X_{(1)}+\frac{b}{2})^{2}-(X_{(n)}-\frac{b}{2})^{2}}{(X_{(1)}+\frac{b}{2})-(X_{(n)}-\frac{b}{2})}=\frac{X_{(1)}+X_{(n)}}{2}.
\end{align*}
\end{example}

\begin{lem}
Suppose that $\rho(t)=t^{2}$. 
\begin{enumerate}
\item If $\delta(X)$ is location equivariant, then its bias $b$ is constant
and $\delta(X)-b$ has smaller risk. 
\item The minimal risk equivariant estimator is unbiased. 
\item If UMVUE exists and is location equivariant, then the MREE and UMVUE
coincide. 
\end{enumerate}
\end{lem}


\subsection{Location-scale equivariance}

We say $(P_{\theta,\tau}:\theta\in\mathbb{R},\tau>0)$ is a location-scale
family if $X\sim P_{\theta,\tau}\iff\frac{X-\theta}{\tau}\sim P_{0,1}$.
If densities exist, then
\[
P_{\theta,\tau}(x)=\frac{1}{\tau}f\left(\frac{x-\theta}{\tau}\right)
\]
where $f(x)=P_{0,1}(x)$. An estimator $\widehat{\theta}(X)=\widehat{\theta}(X_{1},\ldots,X_{n})$
is said to be location-scale equivariant if 
\[
\widehat{\theta}(aX+b)=a\widehat{\theta}(X)+b
\]
for all $a>0$ and $b\in\mathbb{R}$. We say the loss function $L$
is location-scale invariant if
\[
L(a\widehat{\theta}+b,a\theta+b,a\tau)=L(\widehat{\theta},\theta,\tau).
\]
Similar to location invariance, we have $L(\widehat{\theta},\theta,\tau)=L(\widehat{\theta}-\theta,0,\tau)=L(\frac{\widehat{\theta}-\theta}{\tau},0,1)=:\rho(\frac{\widehat{\theta}-\theta}{\tau})$,
so the loss function can be written as a function of $\frac{\widehat{\theta}-\theta}{\tau}$.
Also, the risk of a location-scale equivariant estimator is constant,
since
\[
\mathbb{E}_{\theta,\tau}\rho\left(\frac{\widehat{\theta}(X)-\theta}{\tau}\right)=\mathbb{E}_{\theta,\tau}\rho\left(\widehat{\theta}\left(\frac{X-\theta}{\tau}\right)\right)=\mathbb{E}_{0,1}\rho\left(\widehat{\theta}\left(X\right)\right).
\]
How do we find a minimal risk location-scale eqivariant estimator
$\delta(X)$? Again, we start with pilot estimators, but this time
we need two. Suppose that $\delta_{0}(X)$ and $\delta_{1}(X)$ are
pilot estimators satisfying
\begin{align*}
\delta_{0}(aX+b) & =a\delta_{0}(X)+b,\\
\delta_{1}(aX+b) & =a\delta_{1}(X).
\end{align*}
Then letting $u(X):=\frac{\delta(X)-\delta_{0}(X)}{\delta_{1}(X)}$,
we have
\begin{align*}
u(aX+b) & =\frac{\delta(aX+b)-\delta_{0}(aX+b)}{\delta_{1}(aX+b)}\\
 & =\frac{a\delta(X)+b-a\delta_{0}(X)-b}{a\delta_{1}(X)}\\
 & =\frac{\delta(X)-\delta_{0}(X)}{\delta_{1}(X)}.
\end{align*}
Therefore, $u(aX+b)=u(X)$, and we have
\begin{align*}
u(X_{1},\ldots,X_{n}) & =u(X_{1}-X_{n},\ldots,X_{n-1}-X_{n},0)\\
 & =u\left(\frac{X_{1}-X_{n}}{|X_{n-1}-X_{n}|},\ldots,\frac{X_{n-1}-X_{n}}{|X_{n-1}-X_{n}|},0\right).
\end{align*}
Let $Z=\left(\frac{X_{1}-X_{n}}{|X_{n-1}-X_{n}|},\ldots,\frac{X_{n-1}-X_{n}}{|X_{n-1}-X_{n}|}\right)^{T}$
and write $u(X)=w(Z)$. Note that $Z$ is ancillary of $(\theta,\tau)$.
We have 
\[
\delta(X)=\delta_{0}(X)-w(Z)\delta_{1}(X).
\]
The risk is
\begin{align*}
\mathbb{E}_{0,1}\rho(\delta(X)) & =\mathbb{E}_{0,1}\rho(\delta_{0}(X)-w(Z)\delta_{1}(X))\\
 & =\mathbb{E}_{0,1}\left[\mathbb{E}_{0,1}(\rho(\delta_{0}(X)-w(Z)\delta_{1}(X))|Z)\right],
\end{align*}
so to minimize the risk we need to take
\[
w(Z)=\argmin_{a}\mathbb{E}_{0,1}\left(\rho(\delta_{0}(X))-a\delta_{1}(X)|Z\right).
\]
If $\rho(t)=t^{2}$, then
\begin{align*}
w(Z) & =\argmin_{a}\mathbb{E}_{0,1}\left[\left(\delta_{0}(X)-a\delta_{1}(X)\right)^{2}|Z\right]\\
 & =\argmin_{a}\left[a^{2}\mathbb{E}_{0,1}(\delta_{1}(X)^{2}|Z)-2a\mathbb{E}_{0,1}(\delta_{0}(X)\delta_{1}(X)|Z)\right]\\
 & =\frac{\mathbb{E}_{0,1}(\delta_{0}(X)\delta_{1}(X)|Z)}{\mathbb{E}_{0,1}(\delta_{1}(X)^{2}|Z)}.
\end{align*}
Therefore, 
\[
\delta(X)=\delta_{0}(X)-\delta_{1}(X)\frac{\mathbb{E}_{0,1}(\delta_{0}(X)\delta_{1}(X)|Z)}{\mathbb{E}_{0,1}(\delta_{1}(X)^{2}|Z)}.
\]

\begin{example}
Let $X_{1},\ldots,X_{n}\overset{\text{i.i.d.}}{\sim}E(\theta,b)$.
Recall that $P_{\theta,b}(x)=\frac{1}{b}e^{-(x-\theta)/b}1_{\{x\geq\theta\}}$.
Then $(\delta_{0}(X),\delta_{1}(X))=(X_{(1)},\sum_{i=1}^{n}(X_{i}-X_{(1)}))$
is complete and sufficient for $(\theta,b)$. By Basu's theorem, $(\delta_{0}(X),\delta_{1}(X))$
is independent of $Z$. Also, Basu's theorem implies that $X_{(1)}$
and $\sum_{i=1}^{n}(X_{i}-X_{(1)})$ are independent because $X_{(1)}$
is complete and sufficient for $\theta$ but $\sum_{i=1}^{n}(X_{i}-X_{(1)})$
is ancillary for $\theta$. Therefore, 
\[
\frac{\mathbb{E}_{0,1}(\delta_{0}(X)\delta_{1}(X)|Z)}{\mathbb{E}_{0,1}(\delta_{1}(X)^{2}|Z)}=\frac{\mathbb{E}_{0,1}(\delta_{0}(X))\mathbb{E}(\delta_{1}(X))}{\mathbb{E}_{0,1}(\delta_{1}(X)^{2})}.
\]
Recall from a previous example that $\mathbb{E}(X_{(1)})=\frac{1}{n}$.
The first moments of $\delta_{0}(X)$ and $\delta_{1}(X)$ can be
easily derived. For the second moment of $\delta_{1}(X)$, we will
derive the full distribution of $\delta_{1}(X)$. We have
\begin{align*}
f_{X_{(1)},\ldots,X_{(n)}}(x_{1},\ldots,x_{n}) & =n!f_{X_{1},\ldots,X_{n}}(x_{1},\ldots,x_{n})1_{\{0\leq x_{1}\leq x_{2}\leq\cdots\leq x_{n}\}}\\
 & =n!e^{-\sum_{i=1}^{n}x_{i}}1_{\{0\leq x_{1}\leq x_{2}\leq\cdots\leq x_{n}\}}.
\end{align*}
Let us now make the change of variable $Y_{1}=nX_{(1)}$, $Y_{2}=(n-1)(X_{(2)}-X_{(1)})$,
$Y_{3}=(n-2)(X_{(3)}-X_{(2)}),\ldots,Y_{n}=X_{(n)}-X_{(n-1)}$. Then
\[
\left|\frac{\partial Y}{\partial X_{(1,\ldots,n)}}\right|=n!
\]
so we have
\[
f_{Y_{1},\ldots,Y_{n}}(y_{1},\ldots,y_{n})=f_{X_{(1)},\ldots,X_{(n)}}(x_{1},\ldots,x_{n})\left|\frac{\partial X}{\partial Y}\right|=e^{-\sum_{i=1}^{n}x_{i}}1_{\{0\leq x_{1}\leq x_{2}\leq\cdots\leq x_{n}\}}.
\]
Since $\sum_{i=1}^{n}X_{i}=\sum_{i=1}^{n}Y_{i}$, it follows that
\[
f_{Y_{1},\ldots,Y_{n}}(y_{1},\ldots,y_{n})=e^{-\sum_{i=1}^{n}y_{i}}1_{\{y_{1}\geq0,\ldots,y_{n}\geq0\}}=\prod_{i=1}^{n}e^{-y_{i}}1_{\{y_{i}\geq0\}}.
\]
This shows that $Y_{1},\ldots,Y_{n}\overset{\text{i.i.d.}}{\sim}E(0,1)$.
Then 
\[
\delta_{1}(X)=\sum_{i=1}^{n}(X_{i}-X_{(1)})=\sum_{i=1}^{n}X_{i}-nX_{(1)}=\sum_{i=2}^{n}Y_{i}.
\]
Therefore, $\mathbb{E}_{0,1}(\delta_{1}(X))=n-1$ and 
\begin{align*}
\mathbb{E}_{0,1}(\delta_{1}(X)^{2}) & =\operatorname{Var}_{0,1}(\delta_{1}(X))+(\mathbb{E}_{0,1}\delta_{1}(X))^{2}\\
 & =n-1+(n-1)^{2}\\
 & =n(n-1).
\end{align*}
It follows that
\[
\frac{\mathbb{E}_{0,1}(\delta_{0}(X))\mathbb{E}(\delta_{1}(X))}{\mathbb{E}_{0,1}(\delta_{1}(X)^{2})}=\frac{\frac{1}{n}(n-1)}{n(n-1)}=\frac{1}{n^{2}},
\]
and we have the MREE
\[
\delta(X)=X_{(1)}-\frac{1}{n^{2}}\sum_{i=1}^{n}(X_{i}-X_{(1)}).
\]
\end{example}

The following should be obvious. 
\begin{lem}
If $\delta(X)$ is the minimal risk location equivariant estimator
for $\theta$ when $\tau$ is known, $\delta(X)$ does not depend
on $\tau$, and $\delta(X)$ is also location-scale equivariant, then
$\delta(X)$ is the minimal risk location-scale equivariant estimator
for $\theta$ when $\tau$ is unknown. 
\end{lem}

\begin{example}
Suppose that $X_{1},\ldots,X_{n}\overset{\text{i.i.d.}}{\sim}N(\theta,\sigma^{2})$.
Then $\delta(X)=\overline{X}$ is the minimal risk location equivariant
estimator for $\theta$ when $\sigma^{2}$ is known, $\delta(X)$
does not depend on $\sigma^{2}$, and $\delta(X)$ is location-scale
equivariant, so $\overline{X}$ is the minimal risk location-scale
equivariant estimator for $\theta$ when $\sigma^{2}$ is unknown. 
\end{example}


\section{Likelihood Method and Asymptotic Theory}

\subsection{Consistency of MLE}

Suppose that $X_{1},\ldots,X_{n}\overset{\text{i.i.d.}}{\sim}P_{\theta^{\ast}}$.
The maximum likelihood estimator is defined to be
\[
\widehat{\theta}=\argmax_{\theta}\sum_{i=1}^{n}\log p_{\theta}(X_{i}).
\]
There are a few questions to consider: 
\begin{itemize}
\item Is $\widehat{\theta}$ consistent? That is, can we show that for every
$\epsilon>0$, we have
\[
\mathbb{P}_{\theta^{\ast}}\left(\Vert\widehat{\theta}-\theta^{\ast}\Vert>\epsilon\right)\overset{n\to\infty}{\longrightarrow}0?
\]
\item What is the rate of convergence? 
\item What is the asymptotic distribution of $\widehat{\theta}$? 
\end{itemize}
%
For any two distributions $P$ and $Q$ with densities $p$ and $q$,
we define the Kullback--Leibler divergence $D(P\Vert Q)$ to be
\[
D(P\Vert Q)=\int p\log\frac{p}{q}.
\]

\begin{lem}
$D(P\Vert Q)\geq0$. 
\end{lem}

\begin{proof}
We have
\[
D(P\Vert Q)=\int p\log\frac{p}{q}=\int qf(\frac{p}{q}),
\]
where $f(t)=t\log t$. Since $f$ is convex, by Jensen's inequality,
we have
\begin{align*}
\int qf(\frac{p}{q}) & =\mathbb{E}_{X\sim Q}f\left(\frac{p(X)}{q(X)}\right)\\
 & \geq f\left(\mathbb{E}_{X\sim Q}\frac{p(X)}{q(X)}\right)\\
 & =f(1)\\
 & =0.
\end{align*}
\end{proof}
%
In fact, we have a strong result.
\begin{lem}
$D(P\Vert Q)\geq H^{2}(P,Q)$, where $H^{2}(P,Q)=\frac{1}{2}\int(\sqrt{p}-\sqrt{q})^{2}$
is the squared Hellinger distance, 
\end{lem}

We note that 
\[
\widehat{\theta}=\argmin_{\theta}\frac{1}{n}\sum_{i=1}^{n}\log\frac{p_{\theta^{\ast}}(X_{i})}{p_{\theta}(X_{i})}.
\]
By the law of large numbers, for each fixed $\theta$, 
\[
\frac{1}{n}\sum_{i=1}^{n}\log\frac{p_{\theta^{\ast}}(X_{i})}{p_{\theta}(X_{i})}\overset{n\to\infty}{\longrightarrow}\int p_{\theta^{\ast}}\log\frac{p_{\theta^{\ast}}}{p_{\theta}}=D(P_{\theta^{\ast}}\Vert P_{\theta}).
\]
What is the minimizer of this limiting quantity? Of course, $\theta=\theta^{\ast}$
minimizes the expression, but there could be other values of $\theta$
that minimize the expression if we have $P_{\theta_{1}}=P_{\theta_{2}}$
for some $\theta_{1}\neq\theta_{2}$. To avoid this, we need to assume
$\theta_{1}\neq\theta_{2}\implies P_{\theta_{1}}\neq P_{\theta_{2}}$.
We will make the stronger assumption that for all $\epsilon>0$, there
exists $\delta>0$ such that
\[
\Vert\theta_{1}-\theta_{2}\Vert>\epsilon\implies D(P_{\theta_{1}}|P_{\theta_{2}})>\delta.
\]
This is called identifiability. Then given $\epsilon>0$, we have
\begin{align*}
\mathbb{P}_{\theta^{\ast}}\left(\Vert\widehat{\theta}-\theta^{\ast}\Vert>\epsilon\right) & \leq\mathbb{P}_{\theta^{\ast}}\left(D(P_{\theta^{\ast}}|P_{\widehat{\theta}})>\delta\right).
\end{align*}
We have $D(P_{\theta^{\ast}}|P_{\widehat{\theta}})=\int p_{\theta^{\ast}}\log p_{\theta^{\ast}}-\int p_{\theta^{\ast}}\log p_{\widehat{\theta}}$
and for large $n$, 
\[
\int p_{\theta^{\ast}}\log p_{\theta^{\ast}}\approx\frac{1}{n}\sum_{i=1}^{n}\log p_{\theta^{\ast}}(X_{i})\quad\text{and}\quad\int p_{\theta^{\ast}}\log p_{\widehat{\theta}}\approx\frac{1}{n}\sum_{i=1}^{n}\log p_{\widehat{\theta}}(X_{i}).
\]
The key intuition is that because $\frac{1}{n}\sum_{i=1}^{n}\log p_{\theta^{\ast}}(X_{i})$
is smaller than $\frac{1}{n}\sum_{i=1}^{n}\log p_{\widehat{\theta}}(X_{i})$
by definition of $\widehat{\theta}$, we should expect $\mathbb{P}_{\theta^{\ast}}\left(D(P_{\theta^{\ast}}|P_{\widehat{\theta}})>\delta\right)$
to approach 0 as $n\to\infty$. By the law of large numbers, we have
\[
\frac{1}{n}\sum_{i=1}^{n}\log p_{\theta^{\ast}}(X_{i})\overset{n\to\infty}{\longrightarrow}\int p_{\theta^{\ast}}\log p_{\theta^{\ast}}.
\]
However, we cannot directly apply the law of large numbers to obtain
\[
\frac{1}{n}\sum_{i=1}^{n}\log p_{\widehat{\theta}}(X_{i})\overset{n\to\infty}{\longrightarrow}\int p_{\theta^{\ast}}\log p_{\widehat{\theta}}
\]
because $\log p_{\widehat{\theta}}(X_{1}),\log P_{\widehat{\theta}}(X_{2}),\ldots$
are not i.i.d. Instead, we make use of the uniform law of large numbers
\[
\sup_{\theta}\left|\frac{1}{n}\sum_{i=1}^{n}\log p_{\theta}(X_{i})-\int p_{\theta^{\ast}}\log p_{\theta}\right|\stackrel[n\to\infty]{P_{\theta^{\ast}}}{\longrightarrow}0
\]
which holds under some VC-dimension assumptions. Then we have
\[
\mathbb{P}_{\theta^{\ast}}\left(\sup_{\theta}\left|\frac{1}{n}\sum_{i=1}^{n}\log p_{\theta}(X_{i})-\int p_{\theta^{\ast}}\log p_{\theta}\right|>\frac{\delta}{2}\right)\overset{n\to\infty}{\longrightarrow}0.
\]
We have
\begin{align*}
D(P_{\theta^{\ast}}\Vert P_{\widehat{\theta}}) & =\int p_{\theta^{\ast}}\log p_{\theta^{\ast}}-\int p_{\theta^{\ast}}\log p_{\widehat{\theta}}\\
 & \leq\frac{1}{n}\sum_{i=1}^{n}\log p_{\theta^{\ast}}(X_{i})-\frac{1}{n}\sum_{i=1}^{n}\log p_{\widehat{\theta}}(X_{i})+\left|\frac{1}{n}\sum_{i=1}^{n}\log p_{\theta^{\ast}}(X_{i})-\int p_{\theta^{\ast}}\log p_{\theta^{\ast}}\right|+\left|\frac{1}{n}\sum_{i=1}^{n}\log p_{\widehat{\theta}}(X_{i})-\int p_{\theta^{\ast}}\log p_{\widehat{\theta}}\right|\\
 & \leq\frac{1}{n}\sum_{i=1}^{n}\log p_{\widehat{\theta}}(X_{i})-\frac{1}{n}\sum_{i=1}^{n}\log p_{\widehat{\theta}}(X_{i})+2\sup_{\theta}\left|\frac{1}{n}\sum_{i=1}^{n}\log p_{\theta}(X_{i})-\int p_{\theta^{\ast}}\log p_{\theta}\right|\\
 & =2\sup_{\theta}\left|\frac{1}{n}\sum_{i=1}^{n}\log p_{\theta}(X_{i})-\int p_{\theta^{\ast}}\log p_{\theta}\right|.
\end{align*}
Therefore, 
\[
\mathbb{P}_{\theta^{\ast}}\left(D(P_{\theta^{\ast}}\Vert P_{\widehat{\theta}})>\delta\right)\leq\mathbb{P}_{\theta^{\ast}}\left(2\sup_{\theta}\left|\frac{1}{n}\sum_{i=1}^{n}\log p_{\theta}(X_{i})-\int p_{\theta^{\ast}}\log p_{\theta}\right|>\delta\right)
\]
which converges to 0 as $n\to\infty$. Another way to write the consistency
of the MLE $\widehat{\theta}$ is $\Vert\widehat{\theta}-\theta^{\ast}\Vert=o_{P_{\theta^{\ast}}}(1)$. 

\subsection{Asymptotic normality of MLE}

We will now derive the asymptotic distribution of MLE, under the assumption
that $\Vert\widehat{\theta}-\theta^{\ast}\Vert=o_{P_{\theta^{\ast}}}(1)$
and some other regularity conditions. We will assume in our arguments
that $\theta$ is a scalar, but all arguments can be easily modified
to work with vector parameters. Let us define the score function to
be $S_{\theta}(x)=\frac{\partial}{\partial\theta}\log p_{\theta}(x)$.
We will always assume that the distribution is sufficient regular
to allow interchanging differentiation and integration. Then we have
$\mathbb{E}_{\theta}S_{\theta}(X)=0$ because
\[
\mathbb{E}_{\theta}S_{\theta}(X)=\int p_{\theta}\frac{\partial}{\partial\theta}\log p_{\theta}=\int p_{\theta}\frac{\frac{\partial}{\partial\theta}p_{\theta}}{p_{\theta}}=\int\frac{\partial}{\partial\theta}p_{\theta}=\frac{\partial}{\partial\theta}\int p_{\theta}=0.
\]
We define the Fisher information to be $I_{\theta}=\operatorname{Var}_{\theta}(S_{\theta}(X))=\mathbb{E}_{\theta}\left(S_{\theta}(X)\right)^{2}$.
The claim is that $\sqrt{n}(\widehat{\theta}-\theta^{\ast})\rightsquigarrow N\left(0,I_{\theta^{\ast}}^{-1}\right).$ 
\begin{rem}
If the log-likelihood is twice differentiable, an alternative formula
for the Fisher information is
\[
I_{\theta}=-\mathbb{E}_{\theta}\left(\frac{\partial^{2}}{\partial\theta^{2}}\log p_{\theta}(X)\right).
\]
You can check this easily.
\end{rem}


\subsubsection{The intuition}

We assume that the log-likelihood has a second derivative. Using a
Taylor expansion, we have
\begin{align*}
\frac{1}{n}\sum_{i=1}^{n}\log p_{\theta}(X_{i}) & \approx\frac{1}{n}\sum_{i=1}^{n}\log p_{\theta^{\ast}}(X_{i})+(\theta-\theta^{\ast})\frac{1}{n}\sum_{i=1}^{n}\left.\frac{\partial}{\partial\theta}\log p_{\theta}(X_{i})\right|_{\theta=\theta^{\ast}}+\frac{(\theta-\theta^{\ast})^{2}}{2}\frac{1}{n}\sum_{i=1}^{n}\left.\frac{\partial^{2}}{\partial\theta^{2}}\log p_{\theta}(X_{i})\right|_{\theta=\theta^{\ast}}\\
 & \approx\frac{1}{n}\sum_{i=1}^{n}\log p_{\theta^{\ast}}(X_{i})+(\theta-\theta^{\ast})\frac{1}{n}\sum_{i=1}^{n}S_{\theta^{\ast}}(X_{i})-\frac{1}{2}(\theta-\theta^{\ast})^{2}I_{\theta^{\ast}}=:\widetilde{L}_{n}(\theta).
\end{align*}
Then $\widetilde{L}_{n}(\theta)$ is clearly concave, so it has a
unique maximizer. Define $\widetilde{\theta}=\argmax_{\theta}\widetilde{L}_{n}(\theta)$.
Then setting the derivative of $\widetilde{L}_{n}(\theta)$ to be
0, we have
\[
\widetilde{L}_{n}'(\theta)=\frac{1}{n}\sum_{i=1}^{n}S_{\theta^{\ast}}(X_{i})-(\theta-\theta^{\ast})I_{\theta^{\ast}}=0.
\]
Rearranging gives
\[
\widetilde{\theta}-\theta^{\ast}=\frac{1}{n}\sum_{i=1}^{n}S_{\theta^{\ast}}(X_{i})/I_{\theta^{\ast}}.
\]
Therefore, 
\[
\sqrt{n}\left(\widetilde{\theta}-\theta^{\ast}\right)=\frac{1}{\sqrt{n}}\sum_{i=1}^{n}S_{\theta^{\ast}}(X_{i})/I_{\theta^{\ast}}\rightsquigarrow N\left(0,\operatorname{Var}_{\theta^{\ast}}\left(\frac{S_{\theta^{\ast}}(X_{i})}{I_{\theta^{\ast}}}\right)\right)=N\left(0,I_{\theta^{\ast}}^{-1}\right).
\]
Then since $\widetilde{\theta}\approx\widehat{\theta}$, intuitively
we have also $\sqrt{n}(\widetilde{\theta}-\theta^{\ast})\rightsquigarrow N\left(0,I_{\theta^{\ast}}^{-1}\right).$

\subsubsection{The rigorous argument}

We will now give the rigorous argument which does not assume the existence
of the second derivative of the log-likelihood. Instead of taking
second derivative first and then apply the law of large numbers as
we did above, we will apply the law of large numbers first. The key
idea is that in some cases, the expectation operator can exhibit a
smoothing property that allows $L(\theta):=\mathbb{E}_{\theta^{\ast}}\log p_{\theta}(X)$
to be twice differentiable even if the log-likelihood is not. 
\begin{thm}
[Pollard] Define $L_{n}(\theta)=\frac{1}{n}\sum_{i=1}^{n}\log p_{\theta}(X_{i})$
and define the empirical process operator $\nu_{n}$ by
\[
\nu_{n}f=\frac{1}{\sqrt{n}}\sum_{i=1}^{n}(f(X_{i})-\mathbb{E}_{\theta^{\ast}}f(X)).
\]
Assume that the following conditions hold:
\begin{enumerate}
\item $\Vert\widehat{\theta}-\theta^{\ast}\Vert=o_{P_{\theta^{\ast}}}(1)$.\textup{ }
\item \textup{For every sequence $(U_{n})$ of shrinking neighborhoods of
0, we have 
\[
\sup_{t\in U_{n}}\frac{|\nu_{n}r(\cdot,t)|}{1+\sqrt{n}|t|}\overset{P_{\theta^{\ast}}}{\longrightarrow}0
\]
where $r$ is the remainder term in the first order expansion} $\log p_{\theta^{\ast}+t}(X)=\log p_{\theta^{\ast}}(X)+tS_{\theta^{\ast}}(X)+|t|r(X,t)$.
This condition is known as stochastic differentiability. 
\item $L(\theta)$ has the second order expansion $L(\theta^{\ast}+t)=L(\theta^{\ast})-\frac{t^{2}}{2}I_{\theta^{\ast}}+o(t^{2})$
where $I_{\theta^{\ast}}=\mathbb{E}\left(S_{\theta^{\ast}}(X)\right)^{2}$. 
\end{enumerate}
Then $\sqrt{n}(\widehat{\theta}-\theta^{\ast})\rightsquigarrow N\left(0,I_{\theta^{\ast}}^{-1}\right).$
\end{thm}


\paragraph*{Step 1: Derive a quadratic expansion for $L_{n}(\theta)$}

We have
\begin{align*}
L_{n}(\theta^{\ast}+t) & =L(\theta^{\ast}+t)+\frac{1}{\sqrt{n}}\nu_{n}\log p_{\theta^{\ast}+t}\\
 & =L(\theta^{\ast})-\frac{t^{2}}{2}I_{\theta^{\ast}}+o(t^{2})+\frac{1}{\sqrt{n}}\nu_{n}\left(\log p_{\theta^{\ast}}+tS_{\theta^{\ast}}+|t|r(\cdot,t)\right)\\
 & =\left(L(\theta^{\ast})+\frac{1}{\sqrt{n}}\nu_{n}\log p_{\theta^{\ast}}\right)+\frac{t}{\sqrt{n}}\nu_{n}S_{\theta^{\ast}}-\frac{t^{2}}{2}I_{\theta^{\ast}}+o(t^{2})+\frac{|t|}{\sqrt{n}}\nu_{n}r(\cdot,t)\\
 & =L_{n}(\theta^{\ast})+\frac{t}{n}\sum_{i=1}^{n}S_{\theta^{\ast}}(X_{i})-\frac{t^{2}}{2}I_{\theta^{\ast}}+o(t^{2})+\frac{|t|}{\sqrt{n}}\nu_{n}r(\cdot,t).
\end{align*}
This is almost the same expansion as we had from the intuitive argument. 

\paragraph*{Step 2: Show that $\widehat{\theta}$ is $\sqrt{n}$-consistent}

Let $\widehat{t}=\widehat{\theta}-\theta^{\ast}$. Then
\[
L_{n}(\theta^{\ast}+\widehat{t})=L_{n}(\theta^{\ast})+\frac{\widehat{t}}{n}\sum_{i=1}^{n}S_{\theta^{\ast}}(X_{i})-\frac{\widehat{t}^{2}}{2}I_{\theta^{\ast}}+o_{P_{\theta^{\ast}}}(\widehat{t}^{2})+\frac{|\widehat{t}|}{\sqrt{n}}\nu_{n}r(\cdot,\widehat{t}).
\]
Since $|\widehat{t}|=o_{P_{\theta^{\ast}}}(1)$, we can find a sequence
$(U_{n})$ of shrinking neighbourhoods of 0 such that $\widehat{t}\in U_{n}$.
Thus,
\[
\frac{|\nu_{n}r(\cdot,\widehat{t})|}{1+\sqrt{n}|\widehat{t}|}\leq\sup_{t\in U_{n}}\frac{|\nu_{n}r(\cdot,t)|}{1+\sqrt{n}|t|}=o_{P_{\theta^{\ast}}}(1).
\]
Therefore, 
\[
\frac{|\widehat{t}|}{\sqrt{n}}|\nu_{n}r(\cdot,\widehat{t})|\leq\frac{|\widehat{t}|}{\sqrt{n}}(o_{P_{\theta^{\ast}}}(1)+o_{P_{\theta^{\ast}}}(\sqrt{n}|\widehat{t}|))=o_{P_{\theta^{\ast}}}\left(\frac{|\widehat{t}|}{\sqrt{n}}\right)+o_{P_{\theta^{\ast}}}(\widehat{t}^{2}).
\]
By definition of $\widehat{\theta}$, $L_{n}(\theta^{\ast}+\widehat{t})-L_{n}(\theta^{\ast})\geq0$,
so
\[
0\leq\frac{\widehat{t}}{n}\sum_{i=1}^{n}S_{\theta^{\ast}}(X_{i})-\frac{\widehat{t}^{2}}{2}I_{\theta^{\ast}}+o_{P_{\theta^{\ast}}}\left(\frac{|\widehat{t}|}{\sqrt{n}}\right)+o_{P_{\theta^{\ast}}}(\widehat{t}^{2}).
\]
This gives
\begin{align*}
\frac{\widehat{t}^{2}}{2}I_{\theta^{\ast}} & \leq|\widehat{t}|\left|\frac{1}{n}\sum_{i=1}^{n}S_{\theta^{\ast}}(X_{i})\right|+\left|o_{P_{\theta^{\ast}}}(\widehat{t}^{2})\right|+o_{P_{\theta^{\ast}}}\left(\frac{|\widehat{t}|}{\sqrt{n}}\right)
\end{align*}
which implies that
\[
\left(\frac{1}{2}I_{\theta^{\ast}}-o_{P_{\theta^{\ast}}}(1)\right)|\widehat{t}|\leq\left|\frac{1}{n}\sum_{i=1}^{n}S_{\theta^{\ast}}(X_{i})\right|+o_{P_{\theta^{\ast}}}\left(\frac{1}{\sqrt{n}}\right)=O_{P_{\theta^{\ast}}}\left(\frac{1}{\sqrt{n}}\right),
\]
so $|\widehat{t}|=O_{P_{\theta^{\ast}}}\left(\frac{1}{\sqrt{n}}\right)$. 

\paragraph*{Step 3: Show asymptotic normality}

For any $|t|=O_{P_{\theta}^{\ast}}\left(\frac{1}{\sqrt{n}}\right)$,
we have
\begin{align*}
L_{n}(\theta^{\ast}+t) & =L_{n}(\theta^{\ast})+\frac{t}{n}\sum_{i=1}^{n}S_{\theta^{\ast}}(X_{i})-\frac{t^{2}}{2}I_{\theta^{\ast}}+o_{P_{\theta^{\ast}}}\left(\frac{1}{n}\right)\\
 & =L_{n}(\theta^{\ast})-\frac{1}{2}I_{\theta^{\ast}}\left(t-\frac{1}{n}\sum_{i=1}^{n}\frac{S_{\theta^{\ast}}(X_{i})}{I_{\theta^{\ast}}}\right)^{2}+\frac{1}{2}\frac{\left(\frac{1}{n}\sum_{i=1}^{n}S_{\theta^{\ast}}(X_{i})\right)^{2}}{I_{\theta^{\ast}}}+o_{P_{\theta^{\ast}}}\left(\frac{1}{n}\right).
\end{align*}
Let $\widetilde{t}=\sum_{i=1}^{n}S_{\theta^{\ast}}(X_{i})/I_{\theta^{\ast}}$.
Since both $|\widehat{t}|=O_{P_{\theta^{\ast}}}\left(\frac{1}{\sqrt{n}}\right)$
and $|\widetilde{t}|=O_{P_{\theta^{\ast}}}\left(\frac{1}{\sqrt{n}}\right)$
hold, the inequality $L_{n}(\theta^{\ast}+\widehat{t})\geq L_{n}(\theta^{\ast}+\widetilde{t})$
gives
\[
-\frac{1}{2}I_{\theta^{\ast}}\left(\widehat{t}-\frac{1}{n}\sum_{i=1}^{n}\frac{S_{\theta^{\ast}}(X_{i})}{I_{\theta^{\ast}}}\right)^{2}\geq o_{P_{\theta^{\ast}}}\left(\frac{1}{n}\right).
\]
This implies that 
\[
\frac{1}{2}I_{\theta^{\ast}}\left(\widehat{t}-\frac{1}{n}\sum_{i=1}^{n}\frac{S_{\theta^{\ast}}(X_{i})}{I_{\theta^{\ast}}}\right)^{2}=o_{P_{\theta^{\ast}}}\left(\frac{1}{n}\right)\implies\widehat{t}=\frac{1}{n}\sum_{i=1}^{n}\frac{S_{\theta^{\ast}}(X_{i})}{I_{\theta^{\ast}}}+o_{P_{\theta^{\ast}}}\left(\frac{1}{\sqrt{n}}\right).
\]
Therefore, 
\[
\sqrt{n}(\widehat{\theta}-\theta^{\ast})=\frac{1}{\sqrt{n}}\sum_{i=1}^{n}\frac{S_{\theta^{\ast}}(X_{i})}{I_{\theta^{\ast}}}+o_{P_{\theta^{\ast}}}\left(1\right)\rightsquigarrow N(0,I_{\theta^{\ast}}^{-1}).
\]


\subsection{LAN and DQM}

\subsubsection{Local asymptotic normality}

Let us say that a statistical model $(P_{\theta}:\theta\in\Theta)$
satisfies local asymptotic normality (LAN) at $\theta$ if
\[
\log\prod_{i=1}^{n}\frac{p_{\theta+h/\sqrt{n}}(X_{i})}{p_{\theta}(X_{i})}=\frac{h}{\sqrt{n}}\sum_{i=1}^{n}S_{\theta}(X_{i})-\frac{h^{2}}{2}I_{\theta}+o_{P_{\theta}}\left(1\right)
\]
for any $h$ independent of $n$. Let $\widehat{\theta}$ be the maximum
likelihood estimator. Then LAN at $\theta^{\ast}$ along with consistency
and stochastic differentiability together guarantee that $\sqrt{n}(\widehat{\theta}-\theta^{\ast})\rightsquigarrow N(0,I_{\theta^{\ast}}^{-1})$.
%

\subsubsection{Differentiability in quadratic mean}

It is easy to show that the pointwise derivative $\frac{\partial}{\partial\theta}\sqrt{p_{\theta}(x)}$
is given by $\frac{1}{2}\sqrt{p_{\theta}(x)}S_{\theta}(x)$. If we
assume a slightly stronger notion of differentiability of $\sqrt{p_{\theta}}$,
we can show that LAN holds. Let us say that the parametric family
$(P_{\theta}:\theta\in\Theta)$ satisfies differentiability in quadratic
mean (DQM) at $\theta$ if
\[
\int\left(\frac{\sqrt{p_{\theta+t}(x)}-\sqrt{p_{\theta}(x)}}{t}-\frac{1}{2}\sqrt{p_{\theta}(x)}S_{\theta}(x)\right)^{2}\to0\quad\text{as }t\to0.
\]
In other words, DQM holds if $\theta\mapsto\sqrt{p_{\theta}}$ is
Fréchet differentiable in $L^{2}$ at $\theta$.
\begin{thm}
Suppose that DQM holds at $\theta$. Then $\mathbb{E}_{\theta}S_{\theta}(X)=0$
and LAN holds at $\theta$ with $I_{\theta}=\mathbb{E}_{\theta}(S_{\theta}(X))^{2}$. 
\end{thm}

\textbf{Step 1: Show that $\mathbb{E}_{\theta}S_{\theta}(X)=0$.}
Let $h$ be independent of $n$. Taking $t=h/\sqrt{n}$, we see that
DQM holds at $\theta$ if and only if
\[
\int\left(\sqrt{n}\left(\sqrt{p_{\theta+h/\sqrt{n}}}-\sqrt{p_{\theta}}\right)-\frac{1}{2}h\sqrt{p_{\theta}}S_{\theta}\right)^{2}=o(1).
\]
We have
\begin{align*}
\mathbb{E}_{\theta}S_{\theta}(X) & =\int p_{\theta}S_{\theta}=2\int\frac{1}{2}h\sqrt{p_{\theta}}S_{\theta}\sqrt{p_{\theta}}=I_{1}+I_{2},
\end{align*}
where 
\[
I_{1}=2\int\left(\frac{1}{2}h\sqrt{p_{\theta}}S_{\theta}-\sqrt{n}\left(\sqrt{p_{\theta+h/\sqrt{n}}}-\sqrt{p_{\theta}}\right)\right),
\]
\[
I_{2}=2\int\sqrt{n}\left(\sqrt{p_{\theta+h/\sqrt{n}}}-\sqrt{p_{\theta}}\right)\sqrt{p_{\theta}}.
\]
Then by the Cauchy-Schwarz inequality
\[
|I_{1}|\leq2\sqrt{\int\left(\frac{1}{2}h\sqrt{p_{\theta}}S_{\theta}-\sqrt{n}\left(\sqrt{p_{\theta+h/\sqrt{n}}}-\sqrt{p_{\theta}}\right)\right)^{2}}\sqrt{\int p_{\theta}}=o(1)\overset{n\to\infty}{\longrightarrow}0.
\]
Let $f_{n}=\sqrt{p_{\theta+h/\sqrt{n}}}-\sqrt{p_{\theta}}$ and $g_{n}=\frac{1}{2}\frac{h}{\sqrt{n}}\sqrt{p_{\theta}}S_{\theta}$.
Then
\begin{align*}
|I_{2}| & =2\sqrt{n}\left|\int\sqrt{p_{\theta+h/\sqrt{n}}p_{\theta}}-1\right|\\
 & =\sqrt{n}H^{2}(P_{\theta+h.\sqrt{n}},P_{\theta})\\
 & =\sqrt{n}\left|\int f_{n}^{2}\right|\\
 & =\sqrt{n}\int(f_{n}-g_{n}+g_{n})^{2}\\
 & =\sqrt{n}\int g_{n}^{2}+\sqrt{n}\int(f_{n}-g_{n})^{2}+2\sqrt{n}\int(f_{n}-g_{n})g_{n}.
\end{align*}
We have
\[
\sqrt{n}\int g_{n}^{2}=\sqrt{n}\int\frac{1}{4}\frac{h^{2}}{n}p_{\theta}S_{\theta}^{2}=\frac{h^{2}}{4n}I_{\theta}\overset{n\to\infty}{\longrightarrow}0,
\]
The DQM condition is equivalent to $\int(f_{n}-g_{n})^{2}=o(n^{-1})$,
so we have
\[
\sqrt{n}\int(f_{n}-g_{n})^{2}=o(n^{-1/2})\overset{n\to\infty}{\longrightarrow}0.
\]
Applying the Cauchy-Schwarz inequality, we have
\[
2\sqrt{n}\int(f_{n}-g_{n})g_{n}\leq2\sqrt{n}\sqrt{\int(f_{n}-g_{n})^{2}}\sqrt{\int g_{n}^{2}}\overset{n\to\infty}{\longrightarrow}0.
\]
Hence $I_{2}\overset{n\to\infty}{\longrightarrow}0$ and we conclude
that $\mathbb{E}_{\theta}S_{\theta}(X)=0$. 

\textbf{Step 2: Show that LAN holds at $\theta$.} Note that DQM can
equivalently by stated as
\[
\int p_{\theta}\left(\sqrt{\frac{p_{\theta+h/\sqrt{n}}}{p_{\theta}}}-1-\frac{1}{2}\frac{h}{\sqrt{n}}S_{\theta}\right)^{2}=o(n^{-1}).
\]
Let $W_{i}=\sqrt{\frac{p_{\theta+h/\sqrt{n}}(X_{i})}{p_{\theta}(X_{i})}}-1$.
Then 
\begin{align*}
\log\prod_{i=1}^{n}\frac{p_{\theta+h/\sqrt{n}}(X_{i})}{p_{\theta}(X_{i})} & =2\log\prod_{i=1}^{n}\left(\sqrt{\frac{p_{\theta+h/\sqrt{n}}(X_{i})}{p_{\theta}(X_{i})}}-1+1\right)=2\sum_{i=1}^{n}\log(W_{i}+1).
\end{align*}
Then using the Taylor expansion $\log(x+1)=x-\frac{x^{2}}{2}+o(1)$,
we have
\[
2\sum_{i=1}^{n}\log(W_{i}+1)=2\sum_{i=1}^{n}W_{i}-\sum_{i=1}^{n}W_{i}^{2}+o_{P_{\theta}}(1).
\]
We will analyze the convergence of $\sum_{i=1}^{n}W_{i}$ and $\sum_{i=1}^{n}W_{i}^{2}$.

\textbf{Step 2.1: Analysis for $\sum_{i=1}^{n}W_{i}$. }Let $c_{n}=\sqrt{n}\left(\sqrt{p_{\theta+h/\sqrt{n}}}-\sqrt{p_{\theta}}\right)$
and $c=\frac{1}{2}h\sqrt{p_{\theta}}S_{\theta}$. Since $c_{n}=\sqrt{n}f_{n}$
and $c=\sqrt{n}g_{n}$, it follows from $\int(f_{n}-g_{n})^{2}=o(n^{-1})$
that $\int(c_{n}-c)^{2}=o(1)$. Then
\begin{align*}
\mathbb{E}_{\theta}\sum_{i=1}^{n}W_{i} & =n\left(\int\sqrt{p_{\theta}p_{\theta+h/\sqrt{n}}}-1\right)\\
 & =-\frac{n}{2}\int\left(\sqrt{p_{\theta+h/\sqrt{n}}}-\sqrt{p_{\theta}}\right)^{2}\\
 & =-\frac{1}{2}\int c_{n}^{2}\\
 & =-\frac{1}{2}\int(c_{n}-c+c)^{2}\\
 & =-\frac{1}{2}\int c^{2}-\frac{1}{2}\int(c_{n}-c)^{2}-\int(c_{n}-c)c\\
 & =-\frac{1}{2}\int c^{2}-\int(c_{n}-c)c+o(1).
\end{align*}
Then
\[
-\frac{1}{2}\int c^{2}=-\frac{1}{2}\int\frac{1}{4}h^{2}p_{\theta}S_{\theta}^{2}=-\frac{1}{8}h^{2}I_{\theta},
\]
and by the Cauchy-Schwarz inequality we have
\[
\left|\int(c_{n}-c)c\right|\leq\sqrt{\int(c_{n}-c)^{2}}\sqrt{\int c^{2}}=o(1).
\]
Therefore, 
\[
\mathbb{E}_{\theta}\sum_{i=1}^{n}W_{i}=-\frac{1}{8}h^{2}I_{\theta}+o(1).
\]
Since $\mathbb{E}_{\theta}S_{\theta}(X_{i})=0$, it follows that
\[
\mathbb{E}_{\theta}\sum_{i=1}^{n}\left(W_{i}-\frac{1}{2}\frac{h}{\sqrt{n}}S_{\theta}(X_{i})\right)=-\frac{1}{8}h^{2}I_{\theta}+o(1).
\]
Also,
\begin{align*}
\operatorname{Var}_{\theta}\left(\sum_{i=1}^{n}\left(W_{i}-\frac{1}{2}\frac{h}{\sqrt{n}}S_{\theta}(X_{i})\right)\right) & =n\operatorname{Var}_{\theta}\left(W_{1}-\frac{1}{2}\frac{h}{\sqrt{n}}S_{\theta}(X_{1})\right)\\
 & \leq n\mathbb{E}_{\theta}\left(W_{1}-\frac{1}{2}\frac{h}{\sqrt{n}}S_{\theta}(X_{1})\right)^{2}\\
 & =n\int p_{\theta}\left(\sqrt{\frac{p_{\theta+1/\sqrt{n}}}{p_{\theta}}}-1-\frac{1}{2}\frac{h}{\sqrt{n}}S_{\theta}\right)^{2}\\
 & =n\cdot o(n^{-1})\\
 & =o(1).
\end{align*}
The variance converging to 0 implies that $\sum_{i=1}^{n}\left(W_{i}-\frac{1}{2}\frac{h}{\sqrt{n}}S_{\theta}(X_{i})\right)$
converges in probability to its mean, so we have
\[
\sum_{i=1}^{n}W_{i}-\frac{1}{2}\frac{h}{\sqrt{n}}\sum_{i=1}^{n}S_{\theta}(X_{i})=-\frac{1}{8}h^{2}I_{\theta}+o_{P_{\theta}}(1).
\]
Rearranging gives
\[
\sum_{i=1}^{n}W_{i}=\frac{1}{2}\frac{h}{\sqrt{n}}\sum_{i=1}^{n}S_{\theta}(X_{i})-\frac{1}{8}h^{2}I_{\theta}+o_{P_{\theta}}(1).
\]

\textbf{Step 2.2. Analysis for $\sum_{i=1}^{n}W_{i}^{2}$:} We have
\begin{align*}
\sum_{i=1}^{n}W_{i}^{2} & =\sum_{i=1}^{n}\left(W_{i}-\frac{1}{2}\frac{h}{\sqrt{n}}S_{\theta}(X_{i})+\frac{1}{2}\frac{h}{\sqrt{n}}S_{\theta}(X_{i})\right)^{2}\\
 & =\sum_{i=1}^{n}\left(W_{i}-\frac{1}{2}\frac{h}{\sqrt{n}}S_{\theta}(X_{i})\right)^{2}+\sum_{i=1}^{n}\left(\frac{1}{2}\frac{h}{\sqrt{n}}S_{\theta}(X_{i})\right)^{2}+2\sum_{i=1}^{n}\left(W_{i}-\frac{1}{2}\frac{h}{\sqrt{n}}S_{\theta}(X_{i})\right)\frac{1}{2}\frac{h}{\sqrt{n}}S_{\theta}(X_{i}).
\end{align*}
Call the three terms $S_{1},S_{2}$, and $S_{3}$. For the first term,
we have already found that
\[
\mathbb{E}_{\theta}S_{1}=n\mathbb{E}_{\theta}\left(W_{1}-\frac{1}{2}\frac{h}{\sqrt{n}}S_{\theta}(X_{1})\right)^{2}=o(1).
\]
Therefore, $S_{1}=o_{P_{\theta}}(1)$. For the second term, by the
law of large numbers we have
\[
S_{2}=\frac{h^{2}}{4}\frac{1}{n}\sum_{i=1}^{n}S_{\theta}(X_{i})^{2}=\frac{h^{2}}{4}I_{\theta}+o_{P_{\theta}}(1)=o_{P_{\theta}}(1).
\]
For the third term, by the Cauchy-Schwarz inequality we have
\[
\left|S_{3}\right|\leq2\sqrt{S_{1}}\sqrt{S_{2}}=o_{P_{\theta}}(1).
\]
Thus, we have
\[
\sum_{i=1}^{n}W_{i}^{2}=\frac{h^{2}}{4}I_{\theta}+o_{P_{\theta}}(1).
\]

\textbf{Step 2.3: Conclude. }Putting everything together, we have
\begin{align*}
\log\prod_{i=1}^{n}\frac{p_{\theta+h/\sqrt{n}}(X_{i})}{p_{\theta}(X_{i})} & =2\sum_{i=1}^{n}W_{i}-\sum_{i=1}^{n}W_{i}^{2}+o_{P_{\theta}}(1)\\
 & =2\left(\sum_{i=1}^{n}W_{i}\right)-\sum_{i=1}^{n}W_{i}^{2}+o_{P_{\theta}}(1)\\
 & =2\left(\frac{1}{2}\frac{h}{\sqrt{n}}\sum_{i=1}^{n}S_{\theta}(X_{i})-\frac{1}{8}h^{2}I_{\theta}+o_{P_{\theta}}(1)\right)-\left(\frac{h^{2}}{4}I_{\theta}+o_{P_{\theta}}(1)\right)+o_{P_{\theta}}(1)\\
 & =\frac{h}{\sqrt{n}}\sum_{i=1}^{n}S_{\theta}(X_{i})-\frac{1}{2}h^{2}I_{\theta}+o_{P_{\theta}}(1)
\end{align*}
which shows that LAN holds at $\theta$. 

\subsection{Asymptotic efficiency}

\subsubsection{Cramer-Rao lower bound}

We know from previous discussions that the maximum likelihood estimator
is not always optimal (for example, see discussion on James-Stein
estimator.) But what about asymptotically as $n\to\infty$? Maximum
likelihood estimators are asymptotically efficient in the sense that
it achieves the Cramer-Rao lower bound under the regularity conditions
for asymptotical normality to hold. We state and prove the Cramer-Rao
lower bound for 1-dimensional parameter space. The same result holds
in higher dimensional parameter space. 
\begin{thm}
[Cramer-Rao lower bound] Suppose $X\sim P_{\theta}$ has a score
function and Fisher information. If $\widehat{\theta}$ is unbiased,
then $\operatorname{Var}_{\theta}[\widehat{\theta}]\geq I_{\theta}^{-1}$.
\end{thm}

\begin{proof}
We have 
\[
\operatorname{Var}_{\theta}[\widehat{\theta}]I_{\theta}=\int p_{\theta}(\widehat{\theta}-\theta)^{2}\int p_{\theta}\left(\frac{\frac{\partial}{\partial\theta}p_{\theta}}{p_{\theta}}\right)^{2}\geq\left(\int(\widehat{\theta}-\theta)\frac{\partial}{\partial\theta}p_{\theta}\right)^{2}=\left(\frac{\partial}{\partial\theta}\mathbb{E}_{\theta}[\widehat{\theta}]\right)^{2}=1.
\]
\end{proof}
%
Since the Fisher information of $n$ independent observations is the
sum of each observation's Fisher information, we get the following. 
\begin{cor}
If $X_{1},\ldots,X_{n}\overset{\text{i.i.d.}}{\sim}P_{\theta}$, then
for any unbiased estimator $\widehat{\theta}$, we have $\operatorname{Var}_{\theta}[\widehat{\theta}]\geq(nI_{\theta})^{-1}$. 
\end{cor}


\subsubsection{Hodges' estimator}

Although the maximum likelihood estimator achieves the Cramer-Rao
lower bound, it is possible to find estimators that have lower asymptotic
variance at at least one point. Suppose that $X_{1},\ldots,X_{n}\overset{\text{i.i.d.}}{\sim}N(\theta,1)$
and consider Hodges' estimator, given by
\[
\widehat{\theta}_{H}=\begin{cases}
\overline{X} & |\overline{X}|>n^{-1/4},\\
0 & |\overline{X}|\leq n^{-1/4}.
\end{cases}
\]

\begin{prop}
It holds that
\[
\sqrt{n}(\widehat{\theta}_{H}-\theta)\rightsquigarrow\begin{cases}
N(0,1) & \text{if }\theta\neq0,\\
0 & \text{if }\theta=0.
\end{cases}
\]
\end{prop}

\begin{proof}
Since $X_{1},\ldots,X_{n}\overset{\text{i.i.d.}}{\sim}N(\theta,1)$,
we have $\overline{X}=\theta+\frac{1}{\sqrt{n}}Z$ where $Z\sim N(0,1)$.
Then we can write 
\begin{align*}
\widehat{\theta}_{H} & =\overline{X}1(|\overline{X}|>n^{-1/4})\\
 & =(\theta+\frac{1}{\sqrt{n}}Z)1(|\theta+\frac{1}{\sqrt{n}}Z|>n^{-1/4})\\
 & =\theta-\theta1(|\theta+\frac{1}{\sqrt{n}}Z|\leq n^{-1/4})+\frac{1}{\sqrt{n}}Z1(|\theta+\frac{1}{\sqrt{n}}Z|>n^{-1/4}).
\end{align*}
This gives
\[
\sqrt{n}(\widehat{\theta}_{H}-\theta)=-\sqrt{n}\theta1(|\sqrt{n}\theta+Z|\leq n^{1/4})+Z1(|\sqrt{n}\theta+Z|>n^{1/4}).
\]
Consider the two cases:
\begin{itemize}
\item If $\theta\neq0$, then for every $\omega\in\Omega$ there exists
$N(\omega)$ such that $|\sqrt{n}\theta+Z(\omega)|>n^{1/4}$ for all
$n\geq N(\omega)$. Therefore, $1(|\sqrt{n}\theta+Z|\leq n^{1/4})$
converges to 0 almost surely, and we have $-\sqrt{n}\theta1(|\sqrt{n}\theta+Z|\leq n^{1/4})\to0$
since it is eventually 0. The complementary indicator $1(|\sqrt{n}\theta+Z|>n^{1/4})$
converges to 1, so Slutsky's theorem implies $\sqrt{n}(\widehat{\theta}_{H}-\theta)$
converges in probability to $Z\sim N(0,1)$. 
\item If $\theta=0$, then clearly we have $\sqrt{n}(\widehat{\theta}_{H}-\theta)=Z1(|Z|>n^{1/4})\rightsquigarrow0$. 
\end{itemize}
\end{proof}
Thus, asymptotically Hodges' estimator performs as well as the maximum
likelihood estimator when $\theta\neq0$, but better when $\theta=0$.
Such an estimator is called superefficient and the point 0 is called
a superefficient point. 

\subsubsection{Almost everywhere convolution theorem}

In this section, we look at the almost everywhere convolution theorem,
which can be viewed as an asymptotic version of the Cramer-Rao lower
bound with a much stronger conclusion. Suppose that $(P_{\theta}:\theta\in\Theta)$
satisfies DQM at $\theta$, so that we have
\[
\log\prod_{i=1}^{n}\frac{p_{\theta+h/\sqrt{n}}(X_{i})}{p_{\theta}(X_{i})}=\frac{h}{\sqrt{n}}\sum_{i=1}^{n}S_{\theta}(X_{i})-\frac{h^{2}}{2}I_{\theta}+o_{P_{\theta}}\left(1\right)
\]
for all $h$ independent of $n$. Consider the local experiment $(P_{\theta+h/\sqrt{n}}^{n}:h\in\mathbb{R})$
where local refers to the fact that we treat $h$ as the parameter.
Notice that the right hand side looks like a Gaussian location model. 
\begin{defn}
We write $(P_{\theta+h/\sqrt{n}}^{n}:h\in\mathbb{R})\rightsquigarrow(N(h,I_{\theta}^{-1}):h\in\mathbb{R})$
to denote the expansion above.
\end{defn}

The idea of introducing this notation is to emphasize that the duality
between $(N(h,I_{\theta}^{-1}):h\in\mathbb{R})$ and $(P_{\theta+h/\sqrt{n}}^{n}:h\in\mathbb{R})$.
The following result makes this connection precise. 
\begin{thm}
[Hájek–Le Cam convolution theorem] Suppose that $(P_{\theta}:\theta\in\Theta)$
satisfies DQM at $\theta$. Let $T_{n}$ be a statistic of $(P_{\theta+h/\sqrt{n}}^{n}:h\in\mathbb{R})$
that satisfies
\[
\sqrt{n}\left(T_{n}-\left(\theta+\frac{h}{\sqrt{n}}\right)\right)\overset{P_{\theta+h/\sqrt{n}}^{n}}{\rightsquigarrow}L_{\theta,h}
\]
for every $h\in\mathbb{R}$. Then there exists a randomized statistic
$T$ for $(N(h,I_{\theta}^{-1}):h\in\mathbb{R})$ such that $T-h\overset{N(h,I_{\theta}^{-1})}{\sim}L_{\theta,h}$
for every $h\in\mathbb{R}$. 
\end{thm}

A randomized statistic is a statistic that is allowed to depend on
additional sources of randomness which are independent of the data.
To put it simply, whenever we have an estimator $T_{n}$ for $(P_{\theta+h/\sqrt{n}}^{n}:h\in\mathbb{R})$
with an asymptotic distribution, there exists an estimator for $(N(h,I_{\theta}^{-1}):h\in\mathbb{R})$
whose exact distribution is that same distribution. 
\begin{defn}
Suppose that $T$ is a randomized statistic for $(N(h,I_{\theta}^{-1}):h\in\mathbb{R})$.
We say that $T$ is equivariant in law if $T-h\overset{N(h,I_{\theta}^{-1})}{\sim}L_{\theta}$
for every $h\in\mathbb{R}$.
\end{defn}

The following is a property of Gaussian distributions.
\begin{prop}
If $T$ is a randomized statistic for $(N(h,I_{\theta}^{-1}):h\in\mathbb{R})$
which is equivariant in law satisfying $T-h\overset{N(h,I_{\theta}^{-1})}{\sim}L_{\theta}$
for all $h$, then we have
\[
L_{\theta}=N(0,I_{\theta}^{-1})\ast M_{\theta}
\]
for some distribution $M_{\theta}$. 
\end{prop}

The above observation yields the following. 
\begin{thm}
[convolution theorem] Suppose that $(P_{\theta}:\theta\in\Theta)$
satisfies DQM at $\theta$. Let $T_{n}$ be a statistic of $(P_{\theta+h/\sqrt{n}}^{n}:h\in\mathbb{R})$
that satisfies
\[
\sqrt{n}\left(T_{n}-\left(\theta+\frac{h}{\sqrt{n}}\right)\right)\overset{P_{\theta+h/\sqrt{n}}^{n}}{\rightsquigarrow}L_{\theta}
\]
 for every $h\in\mathbb{R}$. Then $L_{\theta}=N(0,I_{\theta}^{-1})\ast M_{\theta}$
for some distribution $M_{\theta}$. 
\end{thm}

The condition of the above theorem is quite strong, but it can be
replaced by another condition that is much easier to verify.
\begin{lem}
Suppose that $(P_{\theta}:\theta\in\Theta)$ satisfies DQM at all
$\theta\in\Theta\subset\mathbb{R}$ and $T_{n}$ is a statistic satisfying
$\sqrt{n}(T_{n}-\theta)\overset{}{\rightsquigarrow}L_{\theta}$ for
every $\theta\in\Theta$. Then there exists a subsequence $(n_{k})$
such that for Lebesgue almost every $(\theta,h)$, along this subsequence
we have
\[
\sqrt{n}\left(T_{n}-\left(\theta+\frac{h}{\sqrt{n}}\right)\right)\overset{P_{\theta+h/\sqrt{n}}^{n}}{\rightsquigarrow}L_{\theta}.
\]
\end{lem}

This gives the following theorem.
\begin{thm}
[almost everywhere convolution theorem] Suppose that $(P_{\theta}:\theta\in\Theta)$
satisfies DQM at all $\theta\in\Theta\subset\mathbb{R}$ and $T_{n}$
is a statistic satisfying $\sqrt{n}(T_{n}-\theta)\overset{}{\rightsquigarrow}L_{\theta}$
for every $\theta\in\Theta$. Then for Lebesgue almost every $\theta$,
\[
L_{\theta}=N(0,I_{\theta}^{-1})\ast M_{\theta}
\]
for some distribution $M_{\theta}$. 
\end{thm}

In particular, this implies that the set of superefficient points
must have measure zero. 
\end{document}
